\documentclass[11pt,a4paper]{article}

% ── Encoding e lingua ──
\usepackage[utf8]{inputenc}
\usepackage[T1]{fontenc}
\usepackage[italian]{babel}

% ── Layout ──
\usepackage[margin=2.5cm]{geometry}
\usepackage{titlesec}
\usepackage{fancyhdr}
\usepackage{enumitem}

% ── Matematica ──
\usepackage{amsmath, amssymb}

% ── Tabelle ──
\usepackage{booktabs}
\usepackage{longtable}
\usepackage{array}
\usepackage{tabularx}

% ── Grafica ──
\usepackage{graphicx}

% ── Codice ──
\usepackage{listings}
\usepackage{xcolor}

\definecolor{codebg}{rgb}{0.95,0.95,0.95}
\definecolor{codecomment}{rgb}{0.4,0.4,0.4}
\definecolor{codekeyword}{rgb}{0.0,0.0,0.6}
\definecolor{codestring}{rgb}{0.6,0.1,0.1}
\definecolor{codevar}{rgb}{0.1,0.5,0.1}

\lstdefinestyle{sparql}{
  backgroundcolor=\color{codebg},
  basicstyle=\ttfamily\small,
  keywordstyle=\color{codekeyword}\bfseries,
  commentstyle=\color{codecomment}\itshape,
  stringstyle=\color{codestring},
  identifierstyle=\color{codevar},
  breaklines=true,
  frame=single,
  rulecolor=\color{codecomment},
  numbers=none,
  showstringspaces=false,
  tabsize=2,
  captionpos=b,
  morekeywords={PREFIX,SELECT,DISTINCT,WHERE,FILTER,CONSTRUCT,OPTIONAL,
    UNION,GROUP,ORDER,BY,LIMIT,OFFSET,AS,COUNT,SUM,AVG,MIN,MAX,
    BIND,STR,IN,NOT,EXISTS,a,vg,owl,rdf,rdfs,xsd}
}

\lstdefinestyle{turtle}{
  backgroundcolor=\color{codebg},
  basicstyle=\ttfamily\small,
  keywordstyle=\color{codekeyword}\bfseries,
  commentstyle=\color{codecomment}\itshape,
  stringstyle=\color{codestring},
  identifierstyle=\color{codevar},
  breaklines=true,
  frame=single,
  rulecolor=\color{codecomment},
  numbers=none,
  showstringspaces=false,
  tabsize=2,
  captionpos=b,
  morekeywords={@prefix,@base,a,rdf,type,rdfs,subClassOf,owl,Ontology,
    Class,ObjectProperty,DatatypeProperty,inverseOf,domain,range}
}

\lstdefinestyle{python}{
  backgroundcolor=\color{codebg},
  basicstyle=\ttfamily\small,
  keywordstyle=\color{codekeyword}\bfseries,
  commentstyle=\color{codecomment}\itshape,
  stringstyle=\color{codestring},
  breaklines=true,
  frame=single,
  rulecolor=\color{codecomment},
  numbers=none,
  showstringspaces=false,
  tabsize=2,
  captionpos=b
}

% ── Riferimenti e bibliografia ──
\usepackage[colorlinks=true,linkcolor=blue,citecolor=blue,urlcolor=blue]{hyperref}
\usepackage{natbib}
\usepackage{tocbibind}

% ── Header/footer ──
\pagestyle{fancy}
\fancyhf{}
\fancyfoot[C]{\thepage}
\renewcommand{\headrulewidth}{0.4pt}
\fancyhead[L]{\small Videogame Semantic Search}
\fancyhead[R]{\small Relazione Tecnica}

% ── Informazioni documento ──
\title{{\Huge Videogame Semantic Search}\\[6pt]
       {\large Relazione Tecnica di Progetto}\\[6pt]
       {\normalsize Ontologia OWL 2 per videogiochi: modellazione, reasoning e ricerca semantica}}
\author{{\large Videogame Semantic Search Project}}
\date{\today\\[4pt]{\small Versione 1.0}}

\begin{document}

% ── Frontespizio ──
\maketitle
\thispagestyle{empty}

\begin{center}
  \rule{0.8\textwidth}{0.5pt}\\[8pt]
  \textit{Documento tecnico-accademico prodotto nell'ambito del progetto \\
  di Semantic Web --- Ontologia e Knowledge Graph per Videogiochi}\\[8pt]
  \rule{0.8\textwidth}{0.5pt}
\end{center}

\vspace{1cm}

\begin{abstract}
\noindent
Il presente documento descrive il progetto \textbf{Videogame Semantic Search}, un sistema
di ricerca semantica che integra un'ontologia OWL 2 modellata sul dominio dei videogiochi,
un knowledge graph popolato da Wikidata e un'interfaccia web con grafo interattivo.
L'ontologia, costruita secondo i principi del Semantic Web, definisce 78 classi,
31 object property, 26 datatype property e 15 classi definite (pi\`u 2 classi vincolate
da cardinalit\`a massima) classificate tramite reasoning automatico.
La popolazione da Wikidata (2010--2026) produce oltre 1.1 milioni di triple, interrogabili
in linguaggio naturale tramite un agente SPARQL basato su GPT-4.1-mini.
Il documento \`e centrato sull'ontologia: modellazione dello schema e scelte di design,
espressivit\`a DL e profilo OWL 2, reasoning a due livelli (OWL 2 RL + CONSTRUCT) e
interrogazione SPARQL della conoscenza inferita. Popolamento e applicazione web sono
descritti in coda come contesto.
\end{abstract}

\newpage

% ── Indice ──
\tableofcontents
\newpage

% ── Sezioni ──
\section{Introduzione}
\label{sec:introduzione}

\subsection{Contesto: Semantic Web e Ontologie}

Il Semantic Web, introdotto da Tim Berners-Lee nel 2001, rappresenta un'estensione
del Web tradizionale in cui l'informazione \`e strutturata e dotata di semantica
formale, consentendo alle macchine di interpretare e ragionare sui dati in modo
automatico. Il fondamento teorico di questa visione risiede nel concetto di
\textbf{ontologia}, definita da \citet{gruber1993} come \textless{}una specifica
esplicita e formale di una concettualizzazione condivisa\textgreater{}.

Le ontologie, espresse in linguaggi formali come OWL 2 \citep{w3c-owl2-primer},
permettono di modellare un dominio di conoscenza attraverso:
\begin{itemize}[nosep]
  \item \textbf{Classi} --- categorie di entit\`a (es. \texttt{VideoGame}, \texttt{Developer});
  \item \textbf{Propriet\`a} --- relazioni tra entit\`a (es. \texttt{developedBy}) e attributi (es. \texttt{metacriticScore});
  \item \textbf{Assiomi} --- vincoli logici che arricchiscono la semantica (es. classi definite, disjointness, catene di propriet\`a).
\end{itemize}

Un knowledge graph \`e la concretizzazione pratica di un'ontologia popolata con
dati reali: nodi (istanze) e archi (asserzioni RDF) formano un grafo interrogabile
tramite query SPARQL.

\subsection{Obiettivo del Progetto}

Il progetto \textbf{Videogame Semantic Search} si propone di:
\begin{enumerate}[nosep]
  \item \textbf{Costruire un'ontologia OWL 2 di alta qualit\`a} sul dominio dei videogiochi, coprendo il periodo 2010--2026, con uno schema ricco di assiomi (classi primitive e definite, disjointness, property chain, allineamenti esterni);
  \item \textbf{Popolare l'ontologia da Wikidata} tramite query SPARQL all'endpoint pubblico di Wikidata, producendo un insieme di istanze di oltre 1.1 milioni di triple;
  \item \textbf{Abilitare l'interrogazione in linguaggio naturale} attraverso un agente SPARQL basato su Large Language Model (GPT-4.1-mini) che traduce domande in italiano (o inglese) in query SPARQL eseguibili;
  \item \textbf{Visualizzare i risultati come grafo di conoscenza interattivo}, navigabile con zoom, pan, filtri e dettaglio nodi.
\end{enumerate}

\subsection{Casi d'Uso}

I principali casi d'uso del sistema sono:
\begin{itemize}[nosep]
  \item \textbf{Ricerca semantica} --- interrogare il knowledge graph con domande come \textless{}Quali giochi ha sviluppato FromSoftware?\textgreater{} o \textless{}Top 10 giochi RPG con Metacritic pi\`u alto\textgreater{};
  \item \textbf{Grafo di conoscenza interattivo} --- esplorare visualmente le relazioni tra giochi, sviluppatori, generi, piattaforme e franchise;
  \item \textbf{Reasoning automatico} --- classificare automaticamente i giochi in categorie derivate (MultiplayerGame, HighlyRatedGame, SequelGame, ecc.) sfruttando gli assiomi OWL 2.
\end{itemize}

\subsection{Panoramica della Relazione}

Il presente documento \`e organizzato come segue.
Le Sezioni~\ref{sec:ontologia_overview}--\ref{sec:profilo_dl} costituiscono il cuore
della relazione e analizzano in dettaglio l'ontologia: panoramica, classi, propriet\`a
e chiavi, classi definite, disjointness, allineamento esterno, e infine espressivit\`a
DL e profilo OWL 2.
La Sezione~\ref{sec:reasoning} approfondisce il sistema di reasoning a due livelli con
un esempio completo di inferenza, incluso il problema degli \texttt{owl:sameAs} falsi e
la sua risoluzione; la Sezione~\ref{sec:interrogazione} mostra come interrogare con
SPARQL la conoscenza inferita.
Le Sezioni~\ref{sec:popolamento}, \ref{sec:architettura} e \ref{sec:frontend_backend}
descrivono il contesto: la pipeline di popolamento da Wikidata, l'architettura software
e l'applicazione web.
La Sezione~\ref{sec:conclusione} riassume i risultati e delinea possibili sviluppi futuri.
Le appendici riportano le istanze di test, la sintassi DL delle classi definite e gli
assiomi principali dello schema in Turtle.

% Ontologia
\section{Panoramica dell'Ontologia}
\label{sec:ontologia_overview}

\subsection{Dominio}

L'ontologia modella il dominio dei \textbf{videogiochi} nel periodo 2010--2026,
coprendo le seguenti entit\`a principali:
\begin{itemize}[nosep]
  \item \textbf{VideoGame} --- il prodotto videoludico con i suoi attributi (nome, data rilascio, punteggio Metacritic, paese d'origine);
  \item \textbf{Developer / Publisher} --- gli studi di sviluppo e le case editrici;
  \item \textbf{Genre} --- il genere (Action, RPG, Strategy, \dots{});
  \item \textbf{Platform} --- la piattaforma di esecuzione (PS5, PC, Switch, \dots{});
  \item \textbf{Character} --- i personaggi (giocanti e non);
  \item \textbf{Franchise} --- le serie e i sotto-franchise;
  \item \textbf{Award} --- i premi ricevuti;
  \item \textbf{GameEngine} --- il motore grafico utilizzato;
  \item \textbf{GameMode} --- la modalit\`a di gioco (single-player, multiplayer, co-op, MMO);
  \item \textbf{Review / GameEvent} --- recensioni ed eventi del ciclo di vita.
\end{itemize}

\subsection{Namespace e Separazione Schema/Istanze}

Il namespace dell'ontologia \`e:
\begin{center}
\texttt{http://www.videogame-ontology.org/ontology\#}
\end{center}
abbreviato in \texttt{vg:} in tutte le query SPARQL e nei riferimenti interni.

L'ontologia \`e strutturata in due file distinti che realizzano la classica
separazione schema/istanze:
\begin{itemize}[nosep]
  \item \texttt{videogames.owl} --- \textbf{schema}: contiene tutte le dichiarazioni
    di classi, propriet\`a, assiomi, disjointness e allineamenti esterni. Circa 934 triple;
  \item \texttt{videogames\_pruned\_2020.owl} --- \textbf{insieme di istanze di demo}: contiene le istanze
    (giochi, developer, publisher, generi, piattaforme, \dots{}) per il periodo
    2020--2026. Circa 745.000 triple, 68.700 giochi.
\end{itemize}

\subsection{Fonte Dati: Wikidata}

Tutti i dati provengono da \textbf{Wikidata} \citep{wikidata}, il knowledge graph
libero della Wikimedia Foundation. La popolazione avviene tramite query SPARQL
all'endpoint \url{https://query.wikidata.org/sparql}, mappando le propriet\`a Wikidata
rilevanti (identificatori P31, P136, P400, P577, P444, \dots{}) nelle corrispondenti
propriet\`a dell'ontologia.

\subsection{Pipeline di Popolamento}

La pipeline completa si articola in quattro fasi:

\begin{enumerate}[nosep]
  \item \texttt{populate\_wikidata.py} --- esegue 204 iterazioni (17 anni $\times$ 12 mesi)
    di query SPARQL verso Wikidata, raccogliendo giochi, metadati e relazioni.
    Tempo stimato: $\sim$5 ore con polite delay di 1 secondo tra le richieste.
  \item \textbf{OWL-RL closure via \texttt{owlrl}} --- applica la chiusura deduttiva
    OWL 2 RL per materializzare propriet\`a inverse, catene \texttt{subPropertyOf},
    e inferenze da \texttt{someValuesFrom} e \texttt{intersectionOf}. Da $\sim$1M a $\sim$1.4M triple.
  \item \textbf{Pruning} --- rimuove triple ridondanti o inutilizzate a runtime:
    \texttt{owl:sameAs} riflessivi, \texttt{gameDescription}, \texttt{officialWebsite}.
    Da $\sim$1.4M a $\sim$1.17M triple (completo) o $\sim$745k (demo 2020+).
  \item \texttt{enrich\_owl.py} --- gestisce gli assiomi DL-only (FunctionalProperty sulle
    datatype property e \texttt{hasKey}, cfr.\ Sezione~\ref{sec:reasoning}): li rimuove dalla
    copia di lavoro prima del reasoning e, con l'opzione \texttt{-{}-keep-dl-axioms}, li
    ripristina nel file serializzato.
\end{enumerate}

\subsection{Statistiche Riassuntive}

La Tabella~\ref{tab:ontostats} riassume le principali metriche dello schema dell'ontologia.

\begin{table}[h]
\centering
\caption{Statistiche dello schema dell'ontologia}
\label{tab:ontostats}
\begin{tabular}{lr}
\toprule
\textbf{Metrica}                & \textbf{Valore} \\
\midrule
Classi totali                   & 78              \\
Classi primitive                & 61              \\
Classi definite (equivalentClass) & 15            \\
Classi con vincolo di cardinalit\`a & 2           \\
Object Property                 & 31              \\
Datatype Property               & 26              \\
Property chain axioms           & 3               \\
Propriet\`a transitive          & 3               \\
Propriet\`a simmetriche         & 3               \\
Coppie inverse                  & 12              \\
Blocchi AllDisjointClasses      & 12              \\
Vincoli hasKey                  & 3               \\
Allineamenti equivalentClass    & 4 (Wikidata)    \\
Allineamenti equivalentProperty & 3 (Dublin Core) \\
Triple dello schema totali      & $\sim$ 934      \\
\bottomrule
\end{tabular}
\end{table}

Lo schema \`e espresso in formato RDF/XML; l'Appendice~\ref{sec:app_turtle} ne riporta
gli assiomi principali in Turtle. L'espressivit\`a DL e il profilo OWL 2 sono discussi
nella Sezione~\ref{sec:profilo_dl}, il reasoning nella Sezione~\ref{sec:reasoning}.

\section{Classi dello Schema dell'Ontologia}
\label{sec:tbox_classi}

\subsection{Gerarchia delle Classi}

L'ontologia definisce \textbf{12 classi top-level}, ciascuna corrispondente a
un'entit\`a fondamentale del dominio, pi\`u un ricco albero di sottoclassi che
raggiunge un totale di \textbf{78 classi}: 61 primitive, 15 definite tramite
\texttt{owl:equivalentClass} e 2 caratterizzate da un vincolo di cardinalit\`a massima
(condizione solo necessaria).

Le 12 classi top-level sono:
\begin{enumerate}[nosep]
  \item \texttt{vg:VideoGame} --- il videogioco, classe centrale dell'ontologia;
  \item \texttt{vg:Developer} --- lo studio di sviluppo;
  \item \texttt{vg:Publisher} --- l'editore;
  \item \texttt{vg:Genre} --- il genere videoludico;
  \item \texttt{vg:Platform} --- la piattaforma hardware/software;
  \item \texttt{vg:Character} --- il personaggio (giocante o non);
  \item \texttt{vg:Franchise} --- la serie/franchise;
  \item \texttt{vg:Award} --- il premio;
  \item \texttt{vg:GameEngine} --- il motore grafico;
  \item \texttt{vg:GameMode} --- la modalit\`a di gioco;
  \item \texttt{vg:Review} --- la recensione;
  \item \texttt{vg:GameEvent} --- l'evento del ciclo di vita.
\end{enumerate}

\subsection{Tabella Completa delle Classi Primitive}

La Tabella~\ref{tab:classi} elenca tutte le 61 classi primitive (escluse le 15
definite e le 2 con vincolo di cardinalit\`a, trattate nella Sezione~\ref{sec:classi_definite}), raggruppate per
classe top-level di appartenenza.

\begin{longtable}{llp{5cm}}
\caption{Classi primitive dell'ontologia}\label{tab:classi}\\
\toprule
\textbf{Classe} & \textbf{Parent} & \textbf{Descrizione} \\
\midrule
\endfirsthead
\multicolumn{3}{c}{\tablename\ \thetable{} -- continua dalla pagina precedente} \\
\toprule
\textbf{Classe} & \textbf{Parent} & \textbf{Descrizione} \\
\midrule
\endhead
\bottomrule
\multicolumn{3}{r}{\textit{Continua nella pagina successiva}} \\
\endfoot
\bottomrule
\endlastfoot

% ── VideoGame ──
\multicolumn{3}{c}{\textbf{VideoGame (classe top-level)}} \\
\midrule
VideoGame           & owl:Thing       & Prodotto videoludico \\
ReleasedGame        & VideoGame       & Gioco ufficialmente rilasciato \\
EarlyAccessGame     & VideoGame       & Gioco in accesso anticipato \\
CancelledGame       & VideoGame       & Gioco cancellato prima del rilascio \\
DelistedGame        & VideoGame       & Gioco rimosso dagli store digitali \\
IndieGame           & VideoGame       & Gioco a budget ridotto, sviluppatore indipendente \\
AAAGame             & VideoGame       & Gioco ad alto budget, grande publisher \\
GOTYWinner          & AwardWinningGame& Vincitore del premio Game of the Year \\
\midrule

% ── Developer ──
\multicolumn{3}{c}{\textbf{Developer (classe top-level)}} \\
\midrule
Developer           & owl:Thing       & Studio di sviluppo \\
FirstPartyDeveloper & Developer       & Studio proprietario di un platform holder \\
ThirdPartyDeveloper & Developer       & Studio indipendente \\
IndieDeveloper      & Developer       & Piccolo sviluppatore senza publisher \\
\midrule

% ── Publisher ──
\multicolumn{3}{c}{\textbf{Publisher (classe top-level)}} \\
\midrule
Publisher           & owl:Thing       & Casa editrice \\
FirstPartyPublisher & Publisher       & Divisione editoriale di un platform holder \\
ThirdPartyPublisher & Publisher       & Casa editrice indipendente \\
\midrule

% ── Genre ──
\multicolumn{3}{c}{\textbf{Genre (classe top-level)}} \\
\midrule
Genre               & owl:Thing       & Genere videoludico \\
ActionGenre         & Genre           & Azione, sfide fisiche \\
RPGGenre            & Genre           & Gioco di ruolo, progressione \\
StrategyGenre       & Genre           & Strategia, pensiero tattico \\
PuzzleGenre         & Genre           & Rompicapo, problem-solving \\
SimulationGenre     & Genre           & Simulazione di sistemi reali/fittizi \\
SportsGenre         & Genre           & Sport, competizioni atletiche \\
RacingGenre         & Genre           & Corse, competizione su veicoli \\
AdventureGenre      & Genre           & Esplorazione, narrativa \\
HorrorGenre         & Genre           & Horror, paura, suspense \\
PlatformerGenre     & Genre           & Piattaforme, ostacoli \\
\midrule

% ── Platform ──
\multicolumn{3}{c}{\textbf{Platform (classe top-level)}} \\
\midrule
Platform            & owl:Thing       & Piattaforma di esecuzione \\
ConsolePlatform     & Platform        & Console dedicata \\
PlayStationPlatform & ConsolePlatform & Famiglia Sony PlayStation \\
XboxPlatform        & ConsolePlatform & Famiglia Microsoft Xbox \\
NintendoPlatform    & ConsolePlatform & Famiglia Nintendo \\
PCPlatform          & Platform        & Personal computer \\
WindowsPlatform     & PCPlatform      & Microsoft Windows \\
MacOSPlatform       & PCPlatform      & Apple macOS \\
LinuxPlatform       & PCPlatform      & Linux \\
MobilePlatform      & Platform        & Dispositivo mobile \\
IOSPlatform         & MobilePlatform  & Apple iOS \\
AndroidPlatform     & MobilePlatform  & Android \\
\midrule

% ── Character ──
\multicolumn{3}{c}{\textbf{Character (classe top-level)}} \\
\midrule
Character           & owl:Thing       & Personaggio \\
PlayerCharacter     & Character       & Controllato dal giocatore \\
NonPlayerCharacter  & Character       & Controllato dall'IA \\
\midrule

% ── Franchise ──
\multicolumn{3}{c}{\textbf{Franchise (classe top-level)}} \\
\midrule
Franchise           & owl:Thing       & Serie/franchise \\
SubFranchise        & Franchise       & Spin-off o sotto-serie \\
\midrule

% ── Award ──
\multicolumn{3}{c}{\textbf{Award (classe top-level)}} \\
\midrule
Award               & owl:Thing       & Premio \\
IndustryAward       & Award           & Conferito da ente industriale \\
EditorialAward      & Award           & Conferito da testata editoriale \\
UserAward           & Award           & Determinato da voto popolare \\
\midrule

% ── GameEngine ──
\multicolumn{3}{c}{\textbf{GameEngine (classe top-level)}} \\
\midrule
GameEngine          & owl:Thing       & Motore grafico \\
ProprietaryEngine   & GameEngine      & Proprietario, non pubblico \\
OpenSourceEngine    & GameEngine      & Open-source \\
\midrule

% ── GameMode ──
\multicolumn{3}{c}{\textbf{GameMode (classe top-level)}} \\
\midrule
GameMode            & owl:Thing       & Modalit\`a di gioco \\
SinglePlayerMode    & GameMode        & Giocatore singolo \\
MultiplayerMode     & GameMode        & Multigiocatore \\
CoopMode            & GameMode        & Cooperativa \\
MMOMode             & GameMode        & Massively Multiplayer Online \\
\midrule

% ── Review ──
\multicolumn{3}{c}{\textbf{Review (classe top-level)}} \\
\midrule
Review              & owl:Thing       & Recensione (relazione N-aria gioco--revisore--punteggio) \\
\midrule

% ── GameEvent ──
\multicolumn{3}{c}{\textbf{GameEvent (classe top-level)}} \\
\midrule
GameEvent           & owl:Thing       & Evento del ciclo di vita \\
AnnouncementEvent   & GameEvent       & Annuncio pubblico \\
ReleaseEvent        & GameEvent       & Rilascio ufficiale \\
UpdateEvent         & GameEvent       & Aggiornamento/patch \\
DelistingEvent      & GameEvent       & Rimozione dagli store \\
\bottomrule
\end{longtable}

\subsection{Scelte di Modellazione}

\subsubsection{Partizione di Platform}

La classe \texttt{Platform} \`e stata partizionata in tre rami principali
(\texttt{ConsolePlatform}, \texttt{PCPlatform}, \texttt{MobilePlatform}),
ciascuno ulteriormente specializzato:

\begin{verbatim}
Platform
  +-- ConsolePlatform
  |     +-- PlayStationPlatform
  |     +-- XboxPlatform
  |     +-- NintendoPlatform
  +-- PCPlatform
  |     +-- WindowsPlatform
  |     +-- MacOSPlatform
  |     +-- LinuxPlatform
  +-- MobilePlatform
        +-- IOSPlatform
        +-- AndroidPlatform
\end{verbatim}

Questa struttura consente interrogazioni a diversi livelli di granularit\`a:
\textless{}giochi disponibili su console\textgreater{} (\texttt{ConsolePlatform}),
\textless{}giochi per PlayStation\textgreater{} (\texttt{PlayStationPlatform}),
oppure \textless{}giochi per PC Windows\textgreater{} (\texttt{WindowsPlatform}).
La partizione \`e resa semanticamente rigorosa da 4 blocchi \texttt{AllDisjointClasses}
(cfr.~Sezione~\ref{sec:disjointness}). Ogni piattaforma Wikidata viene classificata
al livello pi\`u specifico tramite le regole di reasoning (cfr.~Sezione~\ref{sec:reasoning}).

\subsubsection{Generi come Classi Aperte}

\texttt{Genre} ha 10 sottoclassi (\texttt{ActionGenre}, \texttt{RPGGenre}, \dots{})
ma \textbf{non} viene dichiarata come esaustiva (\texttt{owl:oneOf}) n\'e come
coperta da un'unione. Questa scelta \`e deliberata: Wikidata pu\`o contenere generi
non previsti dalla tassonomia attuale (es. \textless{}Battle Royale\textgreater{},
\textless{}Roguelike\textgreater{}). Un'ontologia chiusa impedirebbe l'aggiunta
di nuovi generi senza modificare lo schema. Con la modellazione aperta, eventuali
generi Wikidata non mappati restano istanze di \texttt{Genre} senza violare
la semantica.

\subsubsection{GameMode come Classe Top-Level Separata}

\texttt{GameMode} \`e modellata come classe top-level separata (anzich\'e come
attributo di \texttt{VideoGame}) per consentire il pattern \textbf{Value Partition}:
le sue 4 sottoclassi (\texttt{SinglePlayerMode}, \texttt{MultiplayerMode},
\texttt{CoopMode}, \texttt{MMOMode}) sono dichiarate mutuamente disgiunte,
garantendo che un gioco con pi\`u modalit\`a abbia asserzioni \texttt{hasGameMode}
distinte e semanticamente non ambigue. Questa scelta permette inoltre di definire
classi derivate (\texttt{SinglePlayerGame}, \texttt{MultiplayerGame}, \texttt{CoopGame})
tramite \texttt{owl:someValuesFrom} sulle modalit\`a.

\subsubsection{Review e GameEvent: Pattern N-ary Relation e Lifecycle Event}

\texttt{Review} e \texttt{GameEvent} sono state introdotte come nuove classi
top-level per modellare relazioni che richiedono reificazione:
\begin{itemize}[nosep]
  \item \textbf{Review} --- reifica la relazione tra un gioco, un revisore e un punteggio.
    Senza la classe \texttt{Review}, sarebbe impossibile associare un punteggio
    (datatype property) a una specifica combinazione gioco-recensore.
    Segue il pattern \textbf{N-ary Relation} (cfr.~Lez.~14);
  \item \textbf{GameEvent} --- reifica eventi temporali come annunci, rilasci e
    aggiornamenti. Consente di associare date e metadati a eventi specifici del
    ciclo di vita di un gioco, seguendo il pattern \textbf{Lifecycle Event}.
\end{itemize}

\section{Propriet\`a dello Schema dell'Ontologia}
\label{sec:tbox_proprieta}

\subsection{Object Property}

La Tabella~\ref{tab:objprop} elenca tutte le 31 object property definite nell'ontologia.
Per semplicit\`a, si omettono le propriet\`a inverse quando speculari; la colonna
\textbf{Assiomi} indica le caratteristiche OWL 2 aggiuntive.

\begin{longtable}{lllll}
\caption{Object Property}\label{tab:objprop}\\
\toprule
\textbf{Propriet\`a}      & \textbf{Domain} & \textbf{Range}   & \textbf{Assiomi}                 & \textbf{Label} \\
\midrule
\endfirsthead
\multicolumn{5}{c}{\tablename\ \thetable{} -- continua} \\
\toprule
\textbf{Propriet\`a}      & \textbf{Domain} & \textbf{Range}   & \textbf{Assiomi}                 & \textbf{Label} \\
\midrule
\endhead
\bottomrule
\endlastfoot

developedBy        & VideoGame & Developer   & Asym, Irr, inv developerOf            & developed by \\
developerOf        & Developer & VideoGame   & Asym, Irr                             & developer of \\
publishedBy        & VideoGame & Publisher   & Asym, Irr, inv publisherOf            & published by \\
publisherOf        & Publisher & VideoGame   & Asym, Irr                             & publisher of \\
hasGenre           & VideoGame & Genre       & Asym, Irr, inv genreOf                & has genre \\
genreOf            & Genre     & VideoGame   & Asym, Irr                             & genre of \\
availableOn        & VideoGame & Platform    & Irr, inv platformFor                  & available on \\
platformFor        & Platform  & VideoGame   & Asym, Irr                             & platform for \\
hasCharacter       & VideoGame & Character   & Asym, Irr, inv appearsIn              & has character \\
appearsIn          & Character & VideoGame   & Asym, Irr                             & appears in \\
belongsTo          & VideoGame & Franchise   & Asym, Irr, inv includes               & belongs to \\
includes           & Franchise & VideoGame   & Asym, Irr                             & includes \\
wonAward           & VideoGame & Award       & Irr, inv awardWonBy                   & won award \\
awardWonBy         & Award     & VideoGame   & Irr                                   & award won by \\
sequelOf           & VideoGame & VideoGame   & Asym, Irr, Trans, inv hasSequel       & sequel of \\
hasSequel          & VideoGame & VideoGame   & Asym, Irr                             & has sequel \\
hasGameMode        & VideoGame & GameMode    & Irr, inv modeInGame                   & has game mode \\
modeInGame         & GameMode  & VideoGame   & Irr                                   & mode in game \\
madeWith           & VideoGame & GameEngine  & Asym, Irr, inv engineUsedIn           & made with \\
engineUsedIn       & GameEngine& VideoGame   & Asym, Irr                             & engine used in \\
partOfFranchise    & Franchise & Franchise   & Asym, Irr, Trans                     & part of franchise \\
subsidiaryOf       & ---       & ---         & Asym, Irr, Trans                     & subsidiary of \\
sharedFranchiseWith& VideoGame & VideoGame   & Sym, chain: belongsTo $\circ$ includes & shared franchise with \\
sharesDeveloperWith& VideoGame & VideoGame   & Sym, chain: developedBy $\circ$ developerOf & shares developer with \\
sharesPublisherWith& VideoGame & VideoGame   & Sym, chain: publishedBy $\circ$ publisherOf & shares publisher with \\
reviewer           & Review    & ---         & ---                                   & reviewer \\
reviewedGame       & Review    & VideoGame   & Funct, inv hasReview                  & reviewed game \\
hasReview          & VideoGame & Review      & ---                                   & has review \\
reviewOf           & Review    & VideoGame   & ---                                   & review of \\
eventFor           & GameEvent & VideoGame   & inv hasEvent                          & event for \\
hasEvent           & VideoGame & GameEvent   & ---                                   & has event \\
\bottomrule
\end{longtable}

\textit{Legenda}: Asym = AsymmetricProperty, Irr = IrreflexiveProperty,
Sym = SymmetricProperty, Trans = TransitiveProperty, Funct = FunctionalProperty,
inv = inverseOf, chain = propertyChainAxiom.

\subsection{Datatype Property}

La Tabella~\ref{tab:dataprop} elenca tutte le 26 datatype property.

\begin{longtable}{lllll}
\caption{Datatype Property}\label{tab:dataprop}\\
\toprule
\textbf{Propriet\`a}      & \textbf{Domain} & \textbf{Range}   & \textbf{subPropertyOf} & \textbf{Label} \\
\midrule
\endfirsthead
\multicolumn{5}{c}{\tablename\ \thetable{} -- continua} \\
\toprule
\textbf{Propriet\`a}      & \textbf{Domain} & \textbf{Range}   & \textbf{subPropertyOf} & \textbf{Label} \\
\midrule
\endhead
\bottomrule
\endlastfoot

name              & ---         & xsd:string  & ---     & name \\
gameName          & VideoGame   & xsd:string  & name    & game name \\
gameDescription   & VideoGame   & xsd:string  & ---     & game description \\
releaseDate       & VideoGame   & xsd:date    & ---     & release date \\
metacriticScore   & VideoGame   & xsd:integer & ---     & Metacritic score \\
developerName     & Developer   & xsd:string  & name    & developer name \\
publisherName     & Publisher   & xsd:string  & name    & publisher name \\
genreName         & Genre       & xsd:string  & name    & genre name \\
platformName      & Platform    & xsd:string  & name    & platform name \\
characterName     & Character   & xsd:string  & name    & character name \\
franchiseName     & Franchise   & xsd:string  & name    & franchise name \\
awardName         & Award       & xsd:string  & name    & award name \\
awardYear         & Award       & xsd:integer & ---     & award year \\
description       & ---         & xsd:string  & ---     & description \\
officialWebsite   & VideoGame   & xsd:anyURI  & ---     & official website \\
countryOfOrigin   & VideoGame   & xsd:string  & ---     & country of origin \\
engineName        & GameEngine  & xsd:string  & name    & engine name \\
hasSteamAppId     & VideoGame   & xsd:string  & ---     & Steam App ID \\
hasWikidataId     & ---         & xsd:string  & ---     & Wikidata ID \\
hasIGDBId         & VideoGame   & xsd:string  & ---     & IGDB ID \\
reviewScore       & Review      & xsd:integer & ---     & review score \\
reviewDate        & Review      & xsd:date    & ---     & review date \\
reviewerName      & Review      & xsd:string  & ---     & reviewer name \\
eventDate         & GameEvent   & xsd:date    & ---     & event date \\
budgetTier        & VideoGame   & xsd:string  & ---     & budget tier \\
playerCount       & VideoGame   & xsd:integer & ---     & player count \\
\bottomrule
\end{longtable}

\subsection{Property Chain Axiom}

L'ontologia definisce tre \texttt{owl:propertyChainAxiom}, che permettono di
inferire relazioni transitive composte. In notazione DL (Description Logic):

\begin{enumerate}[nosep]
  \item \texttt{sharedFranchiseWith} $\equiv$ \texttt{belongsTo} $\circ$ \texttt{includes}

    Due giochi condividono un franchise se il primo appartiene a un franchise
    e il franchise include il secondo.
  \item \texttt{sharesDeveloperWith} $\equiv$ \texttt{developedBy} $\circ$ \texttt{developerOf}

    Due giochi condividono lo sviluppatore se il primo \`e sviluppato da uno
    studio che ha sviluppato anche il secondo.
  \item \texttt{sharesPublisherWith} $\equiv$ \texttt{publishedBy} $\circ$ \texttt{publisherOf}

    Analogo per l'editore.
\end{enumerate}

Queste catene sono materializzate dal reasoner OWL 2 RL (\texttt{owlrl})
e producono relazioni simmetriche tra giochi dello stesso franchise, studio o editore.
Per evitare esplosioni combinatorie su franchise/studio/editore con migliaia di
giochi, il materialiser mirato in \texttt{enrich\_owl.py} limita il numero di
coppie per gruppo tramite il flag \texttt{--max-group} (default 50).

\subsection{Propriet\`a Transitive}

Due propriet\`a sono dichiarate transitive:

\begin{itemize}[nosep]
  \item \texttt{sequelOf} --- se $A$ \`e sequel di $B$ e $B$ \`e sequel di $C$,
    allora $A$ \`e sequel di $C$. Permette di tracciare catene di sequel
    (es. Dark Souls III $\rightarrow$ Dark Souls II $\rightarrow$ Dark Souls);
  \item \texttt{partOfFranchise} --- per gerarchie di franchise con sotto-franchise.
    Es. \textless{}Mario Kart\textgreater{} $\rightarrow$ \textless{}Mario\textgreater{}.
\end{itemize}

La transitivit\`a di \texttt{sequelOf} \`e particolarmente utile per query del
tipo \textless{}tutti i sequel di Dark Souls\textgreater{} senza dover navigare
manualmente la catena.

\subsection{Coppie Inverse Sistematiche}

L'ontologia definisce 13 coppie inverse, consentendo la navigazione bidirezionale
del grafo. Per ogni propriet\`a ``forward'' esiste la corrispondente inversa:
\begin{itemize}[nosep]
  \item \texttt{developedBy} $\leftrightarrow$ \texttt{developerOf}
  \item \texttt{publishedBy} $\leftrightarrow$ \texttt{publisherOf}
  \item \texttt{hasGenre} $\leftrightarrow$ \texttt{genreOf}
  \item \texttt{availableOn} $\leftrightarrow$ \texttt{platformFor}
  \item \texttt{hasCharacter} $\leftrightarrow$ \texttt{appearsIn}
  \item \texttt{belongsTo} $\leftrightarrow$ \texttt{includes}
  \item \texttt{wonAward} $\leftrightarrow$ \texttt{awardWonBy}
  \item \texttt{sequelOf} $\leftrightarrow$ \texttt{hasSequel}
  \item \texttt{hasGameMode} $\leftrightarrow$ \texttt{modeInGame}
  \item \texttt{madeWith} $\leftrightarrow$ \texttt{engineUsedIn}
  \item \texttt{reviewedGame} $\leftrightarrow$ \texttt{hasReview}
  \item \texttt{eventFor} $\leftrightarrow$ \texttt{hasEvent}
\end{itemize}

La tredicesima coppia (\texttt{hasGenre} $\leftrightarrow$ \texttt{genreOf}) era
gi\`a implicita nello schema originale. La dichiarazione esplicita di \texttt{owl:inverseOf}
in una sola direzione \`e sufficiente: \texttt{owlrl} inferisce automaticamente
la direzione opposta durante la materializzazione. Il validatore dello schema segnala
12 warning \texttt{INVERSE\_NOT\_BIDIRECTIONAL} proprio per questa asimmetria
dichiarativa (cfr.~Sezione~\ref{sec:validazione}).

\section{Classi Definite}
\label{sec:classi_definite}

\subsection{Panoramica}

Le classi definite (\textit{defined classes}) costituiscono uno degli aspetti
pi\`u avanzati dell'ontologia. A differenza delle classi primitive, che richiedono
l'asserzione esplicita \texttt{rdf:type}, le classi definite specificano
condizioni necessarie e sufficienti (\texttt{owl:equivalentClass}) che permettono
a un reasoner di \textbf{inferire automaticamente} l'appartenenza di un'istanza.

L'ontologia definisce \textbf{17 classi definite}, raggruppabili in sei categorie
logiche. La Tabella~\ref{tab:definite} le elenca con la sintassi DL e il tipo
di condizione.

\begin{longtable}{lllp{4cm}}
\caption{Classi definite}\label{tab:definite}\\
\toprule
\textbf{Classe} & \textbf{Tipo} & \textbf{Sintassi DL} & \textbf{Costrutti OWL 2} \\
\midrule
\endfirsthead
\multicolumn{4}{c}{\tablename\ \thetable{} -- continua} \\
\toprule
\textbf{Classe} & \textbf{Tipo} & \textbf{Sintassi DL} & \textbf{Costrutti OWL 2} \\
\midrule
\endhead
\bottomrule
\endlastfoot

% Mode-derived
MultiplayerGame     & S/N & VideoGame $\sqcap$ $\exists$hasGameMode.MultiplayerMode   & intersectionOf, someValuesFrom \\
SinglePlayerGame    & S/N & VideoGame $\sqcap$ $\exists$hasGameMode.SinglePlayerMode  & intersectionOf, someValuesFrom \\
CoopGame            & S/N & VideoGame $\sqcap$ $\exists$hasGameMode.CoopMode           & intersectionOf, someValuesFrom \\
% Platform-derived
PCGame              & S/N & VideoGame $\sqcap$ $\exists$availableOn.PCPlatform         & intersectionOf, someValuesFrom \\
ConsoleGame         & S/N & VideoGame $\sqcap$ $\exists$availableOn.ConsolePlatform    & intersectionOf, someValuesFrom \\
MobileGame          & S/N & VideoGame $\sqcap$ $\exists$availableOn.MobilePlatform     & intersectionOf, someValuesFrom \\
NintendoGame        & S/N & VideoGame $\sqcap$ $\exists$availableOn.NintendoPlatform   & intersectionOf, someValuesFrom \\
PlayStationGame     & S/N & VideoGame $\sqcap$ $\exists$availableOn.PlayStationPlatform& intersectionOf, someValuesFrom \\
XboxGame            & S/N & VideoGame $\sqcap$ $\exists$availableOn.XboxPlatform       & intersectionOf, someValuesFrom \\
% Sequel-derived
SequelGame          & S/N & VideoGame $\sqcap$ $\exists$sequelOf.VideoGame             & intersectionOf, someValuesFrom \\
StandaloneGame      & S/N & VideoGame $\sqcap$ $\neg\exists$sequelOf.VideoGame          & intersectionOf, complementOf \\
% Quality-derived
HighlyRatedGame     & S/N & VideoGame $\sqcap$ $\exists$metacriticScore.$\geq_{85}$    & intersectionOf, someValuesFrom, datatype restriction \\
RecentGame          & S/N & VideoGame $\sqcap$ $\exists$releaseDate.$\geq_{2020}$      & intersectionOf, someValuesFrom, datatype restriction \\
% Preesistenti
AwardWinningGame    & S/N & VideoGame $\sqcap$ $\exists$wonAward.Award                 & intersectionOf, someValuesFrom \\
FranchiseGame       & S/N & VideoGame $\sqcap$ $\exists$belongsTo.Franchise            & intersectionOf, someValuesFrom \\
% Cardinality (solo necessarie)
GameWithSingleDeveloper & N & VideoGame $\sqcap$ $\leq1$developedBy                   & subClassOf, maxCardinality \\
ExclusiveRelease    & N & VideoGame $\sqcap$ $\leq1$availableOn                      & subClassOf, maxCardinality \\
\bottomrule
\end{longtable}

\textit{Legenda}: S/N = condizioni sufficienti e necessarie (equivalentClass);
N = solo necessarie (subClassOf).

\subsection{Analisi per Categoria}

\subsubsection{Mode-derived}

Le classi \texttt{MultiplayerGame}, \texttt{SinglePlayerGame} e \texttt{CoopGame}
classificano i giochi in base alla modalit\`a di gioco. Un gioco viene inferito
come \texttt{MultiplayerGame} se esiste almeno un'asserzione \texttt{hasGameMode}
verso un'entit\`a di tipo \texttt{MultiplayerMode}. Si noti che \texttt{SinglePlayerGame}
e \texttt{MultiplayerGame} \textbf{non} sono dichiarate disgiunte: un gioco pu\`o
supportare entrambe le modalit\`a (es. Call of Duty ha sia campagna single-player
che multiplayer).

\textbf{Esempio}: Elden Ring ha \texttt{hasGameMode SinglePlayerMode} e
\texttt{hasGameMode MultiplayerMode} $\rightarrow$ inferito sia come
\texttt{SinglePlayerGame} che come \texttt{MultiplayerGame}.

\subsubsection{Platform-derived}

Sei classi classificano i giochi per piattaforma a tre livelli di granularit\`a:
macrocategoria (\texttt{PCGame}, \texttt{ConsoleGame}, \texttt{MobileGame}) e
famiglia specifica (\texttt{NintendoGame}, \texttt{PlayStationGame}, \texttt{XboxGame}).
L'inferenza sfrutta la gerarchia \texttt{rdfs:subClassOf} delle piattaforme: se
una piattaforma \`e di tipo \texttt{PlayStationPlatform}, il reasoner inferisce
anche \texttt{ConsolePlatform} e \texttt{Platform}.

\textbf{Esempio}: God of War Ragnar\"ok ha \texttt{availableOn} PS5 (tipo
\texttt{PlayStationPlatform}) $\rightarrow$ inferito \texttt{PlayStationGame} e \texttt{ConsoleGame}.

\subsubsection{Sequel-derived}

\texttt{SequelGame} e \texttt{StandaloneGame} formano una partizione completa
dei giochi basata sulla propriet\`a transitiva \texttt{sequelOf}:
\begin{itemize}[nosep]
  \item \texttt{SequelGame} --- $\exists$\texttt{sequelOf.VideoGame}: il gioco ha
    almeno un predecessore diretto;
  \item \texttt{StandaloneGame} --- $\neg\exists$\texttt{sequelOf.VideoGame}: il
    gioco non ha predecessori. Utilizza \texttt{owl:complementOf}, costrutto
    non supportato da OWL 2 RL, e richiede pertanto il reasoner CONSTRUCT
    (cfr.~Sezione~\ref{sec:reasoning}).
\end{itemize}

\textbf{Esempio}: Elden Ring non ha \texttt{sequelOf} $\rightarrow$ \texttt{StandaloneGame};
Dark Souls III ha \texttt{sequelOf} Dark Souls II $\rightarrow$ \texttt{SequelGame}.

\subsubsection{Quality-derived}

\texttt{HighlyRatedGame} e \texttt{RecentGame} utilizzano \textbf{datatype restriction}
con \texttt{xsd:minInclusive}, un costrutto OWL 2 DL non coperto da OWL 2 RL
(cfr.~Lez.~17--18):

\begin{itemize}[nosep]
  \item \texttt{HighlyRatedGame} $\equiv$ VideoGame $\sqcap$ $\exists$\texttt{metacriticScore.}$\geq_{85}$
    --- giochi con Metacritic $\geq$ 85;
  \item \texttt{RecentGame} $\equiv$ VideoGame $\sqcap$ $\exists$\texttt{releaseDate.}$\geq_{2020-01-01}$
    --- giochi rilasciati dal 2020.
\end{itemize}

\textbf{Esempio}: Elden Ring (metacritic 96) $\rightarrow$ \texttt{HighlyRatedGame}.

\subsubsection{Award e Franchise (preesistenti)}

\texttt{AwardWinningGame} e \texttt{FranchiseGame} sono le due classi definite
originali (V1), preservate nell'estensione V2:
\begin{itemize}[nosep]
  \item \texttt{AwardWinningGame}: giochi con almeno un premio (\texttt{wonAward});
  \item \texttt{FranchiseGame}: giochi appartenenti a una serie (\texttt{belongsTo}).
\end{itemize}

\subsubsection{Cardinality (solo necessarie)}

\texttt{GameWithSingleDeveloper} e \texttt{ExclusiveRelease} sono le uniche due
classi definite con \textbf{sole condizioni necessarie} (\texttt{rdfs:subClassOf}
anzich\'e \texttt{owl:equivalentClass}). Questa scelta \`e deliberata:
\begin{itemize}[nosep]
  \item \texttt{GameWithSingleDeveloper}: $\leq 1$ \texttt{developedBy}.
    Condizione sufficiente sarebbe troppo forte (un gioco con esattamente 1
    developer dovrebbe essere \textit{solo} quello);
  \item \texttt{ExclusiveRelease}: $\leq 1$ \texttt{availableOn}.
    Analogo ragionamento: un gioco esclusiva \textit{ha al massimo} 1 piattaforma,
    ma non tutti i giochi con $\leq$ 1 piattaforma sono necessariamente esclusive.
\end{itemize}

Queste classi forniscono vincoli di integrit\`a piuttosto che meccanismi di
classificazione automatica.

\subsection{Costrutti OWL 2 Utilizzati}

La Tabella~\ref{tab:costrutti} riassume i costrutti OWL 2 impiegati nelle classi
definite, con riferimento alle lezioni del corso.

\begin{table}[h]
\centering
\caption{Costrutti OWL 2 nelle classi definite}
\label{tab:costrutti}
\begin{tabular}{lll}
\toprule
\textbf{Costrutto}          & \textbf{Classi}                                   & \textbf{Rif. Lezione} \\
\midrule
\texttt{owl:someValuesFrom}& MultiplayerGame, PCGame, SequelGame, \dots{} (12) & Lez.~15 (esistenziale) \\
\texttt{owl:intersectionOf}& Tutte le 17 definite                             & Lez.~16 (intersezione) \\
\texttt{owl:complementOf}  & StandaloneGame                                   & Lez.~16 (complemento) \\
\texttt{owl:maxCardinality}& GameWithSingleDeveloper, ExclusiveRelease        & Lez.~17--18 \\
\texttt{owl:withRestrictions} & HighlyRatedGame, RecentGame                   & Lez.~17--18 (datatype) \\
\texttt{xsd:minInclusive}  & HighlyRatedGame, RecentGame                      & Lez.~17--18 \\
\bottomrule
\end{tabular}
\end{table}

\section{Disjointness}
\label{sec:disjointness}

\subsection{Blocchi AllDisjointClasses}

L'ontologia dichiara \textbf{12 blocchi \texttt{owl:AllDisjointClasses}}, elencati
nella Tabella~\ref{tab:disjoint}. Ciascun blocco specifica un insieme di classi
mutuamente disgiunte: nessuna istanza pu\`o appartenere a pi\`u di una classe
nello stesso blocco.

\begin{table}[h]
\centering
\caption{Blocchi AllDisjointClasses}
\label{tab:disjoint}
\begin{tabular}{ll}
\toprule
\textbf{\#} & \textbf{Classi} \\
\midrule
1 & VideoGame, Developer, Publisher, Genre, Platform, Character, Franchise, Award, \\
  & GameEngine, GameMode, Review, GameEvent \\
2 & ConsolePlatform, PCPlatform, MobilePlatform \\
3 & PlayStationPlatform, XboxPlatform, NintendoPlatform \\
4 & WindowsPlatform, MacOSPlatform, LinuxPlatform \\
5 & IOSPlatform, AndroidPlatform \\
6 & IndustryAward, EditorialAward, UserAward \\
7 & PlayerCharacter, NonPlayerCharacter \\
8 & ProprietaryEngine, OpenSourceEngine \\
9 & ReleasedGame, EarlyAccessGame, CancelledGame, DelistedGame \\
10& IndieGame, AAAGame \\
11& SinglePlayerMode, MultiplayerMode, CoopMode, MMOMode \\
12& AnnouncementEvent, ReleaseEvent, UpdateEvent, DelistingEvent \\
\bottomrule
\end{tabular}
\end{table}

\subsection{Importanza delle Disjointness}

Le dichiarazioni di disjointness svolgono due ruoli fondamentali:

\begin{enumerate}[nosep]
  \item \textbf{Rilevazione di inconsistenze nei dati Wikidata}. Se un'entit\`a
    Wikidata venisse erroneamente classificata come, ad esempio, sia
    \texttt{ConsolePlatform} che \texttt{PCPlatform}, un reasoner DL completo
    (HermiT, Pellet) rileverebbe l'inconsistenza. Questo meccanismo funge da
    \textit{sanity check} automatico sulla qualit\`a dei dati;
  \item \textbf{Ragionamento DL}. Le disjointness abilitano inferenze per
    contraddizione. Ad esempio, sapendo che \texttt{PlayerCharacter} e
    \texttt{NonPlayerCharacter} sono disgiunte, se un personaggio viene
    classificato come \texttt{PlayerCharacter}, si pu\`o inferire che
    \textbf{non} \`e \texttt{NonPlayerCharacter}.
\end{enumerate}

\subsection{Nota su SinglePlayerGame e MultiplayerGame}

Si noti che \texttt{SinglePlayerGame} e \texttt{MultiplayerGame} \textbf{non}
compaiono in alcun blocco di disjointness. La scelta \`e deliberata: nel dominio
dei videogiochi, \`e comune che un gioco supporti sia la modalit\`a single-player
che multiplayer (es. Grand Theft Auto V, Elden Ring, Call of Duty). Dichiararle
disgiunte produrrebbe inconsistenze su una porzione significativa delle istanze.

Al contrario, le modalit\`a \texttt{SinglePlayerMode}, \texttt{MultiplayerMode},
\texttt{CoopMode} e \texttt{MMOMode} (blocco \#11) \textbf{sono} disgiunte:
una specifica modalit\`a di gioco \`e esattamente una di queste quattro.

\section{Allineamento Esterno}
\label{sec:allineamento}

\subsection{equivalentClass verso Wikidata}

L'ontologia dichiara 4 allineamenti \texttt{owl:equivalentClass} verso entit\`a
Wikidata, riportati nella Tabella~\ref{tab:align_class}.

\begin{table}[h]
\centering
\caption{Allineamenti equivalentClass verso Wikidata}
\label{tab:align_class}
\begin{tabular}{lll}
\toprule
\textbf{Classe VG} & \textbf{Entit\`a Wikidata} & \textbf{Descrizione Wikidata} \\
\midrule
VideoGame & wd:Q7889      & Video game \\
Developer & wd:Q210167    & Video game developer \\
GameEngine& wd:Q193564    & Game engine \\
Genre     & wd:Q65956376  & Video game genre \\
\bottomrule
\end{tabular}
\end{table}

\subsection{equivalentProperty verso Dublin Core}

Tre datatype property sono allineate con il vocabolario Dublin Core \citep{dublincore}
tramite \texttt{owl:equivalentProperty} (Tabella~\ref{tab:align_prop}).

\begin{table}[h]
\centering
\caption{Allineamenti equivalentProperty verso Dublin Core}
\label{tab:align_prop}
\begin{tabular}{lll}
\toprule
\textbf{Propriet\`a VG} & \textbf{Propriet\`a DC} & \textbf{Label DC} \\
\midrule
releaseDate    & dcterms:issued       & Date Issued \\
gameName       & dcterms:title        & Title \\
gameDescription& dcterms:description   & Description \\
\bottomrule
\end{tabular}
\end{table}

\subsection{Vantaggi dell'Allineamento}

L'allineamento esterno, noto in letteratura come \textbf{ontology alignment}
(cfr.~Lez.~13), offre tre vantaggi principali:

\begin{enumerate}[nosep]
  \item \textbf{Interoperabilit\`a}. Un'applicazione che consuma dati Dublin Core
    pu\`o automaticamente interpretare \texttt{vg:releaseDate} come sinonimo di
    \texttt{dcterms:issued}, senza bisogno di mapping manuali;
  \item \textbf{Linkage automatico}. Le entit\`a Wikidata esterne referenziate
    tra le istanze (es. \texttt{wd:Q12345}) sono automaticamente riconosciute come
    istanze di \texttt{vg:VideoGame} se Wikidata le classifica come Q7889;
  \item \textbf{Riuso di standard consolidati}. Anzich\'e reinventare vocabolari,
    l'ontologia si appoggia a standard ampiamente adottati (Dublin Core per
    metadati bibliografici, Wikidata per identificazione entit\`a), seguendo
    il principio \textit{Linked Data} di riusare URI esistenti.
\end{enumerate}

\subsection{Limitazione Attuale}

L'allineamento \`e dichiarato a livello dello schema ma non viene sfruttato dal
reasoner OWL 2 RL, poich\'e \texttt{owlrl} opera solo sul namespace locale.
Per sfruttare appieno questi assiomi, sarebbe necessario un reasoner DL completo
(HermiT, Pellet) che possa importare anche le definizioni di classe da Wikidata.

\section{Profilo OWL 2 ed Espressivit\`a DL}
\label{sec:profilo_dl}

Questa sezione colloca l'ontologia rispetto alla teoria delle Description Logic
e ai profili di OWL 2, e spiega perch\'e lo schema \`e progettato per OWL 2 DL
mentre il reasoning pratico usa il profilo OWL 2 RL (Sezione~\ref{sec:reasoning}).

\subsection{Espressivit\`a: $\mathcal{SRIN}^{(\mathcal{D})}$}

OWL 2 DL corrisponde alla logica descrittiva $\mathcal{SROIQ}^{(\mathcal{D})}$.
I costrutti effettivamente usati nello schema collocano l'ontologia nel frammento
$\mathcal{SRIN}^{(\mathcal{D})}$. La Tabella~\ref{tab:espressivita} associa a ogni
lettera il costrutto e il punto dello schema che lo richiede.

\begin{table}[h]
\centering
\small
\caption{Espressivit\`a DL dell'ontologia}
\label{tab:espressivita}
\begin{tabularx}{\textwidth}{l >{\raggedright\arraybackslash}p{4.8cm} >{\raggedright\arraybackslash}X}
\toprule
\textbf{Lettera} & \textbf{Costrutto} & \textbf{Uso nello schema} \\
\midrule
$\mathcal{ALC}$ & $\sqcap$, $\exists$, $\lnot$ (e $\forall$, $\sqcup$ per dualit\`a) &
  classi definite \texttt{VideoGame} $\sqcap$ $\exists p.C$; complemento in \texttt{StandaloneGame} \\
$\mathcal{S}$   & $\mathcal{ALC}$ + ruoli transitivi & \texttt{sequelOf}, \texttt{partOfFranchise}, \texttt{subsidiaryOf} \\
$\mathcal{R}$   & inclusioni complesse di ruoli, asimmetria, irriflessivit\`a &
  3 property chain; 13 propriet\`a asimmetriche e 18 irriflessive \\
$\mathcal{I}$   & ruoli inversi & 12 \texttt{owl:inverseOf}; 3 propriet\`a simmetriche ($R \equiv R^-$) \\
$\mathcal{N}$   & restrizioni di cardinalit\`a non qualificate &
  ${\le}1\,\texttt{developedBy}$, ${\le}1\,\texttt{availableOn}$; \texttt{reviewedGame} funzionale \\
$(\mathcal{D})$ & datatype, facet, data property funzionali &
  26 datatype property; \texttt{minInclusive} in \texttt{HighlyRatedGame} e \texttt{RecentGame} \\
\bottomrule
\end{tabularx}
\end{table}

Rispetto a $\mathcal{SROIQ}^{(\mathcal{D})}$ mancano due componenti:
\begin{itemize}[nosep]
  \item $\mathcal{O}$ (nominali): lo schema non usa \texttt{owl:oneOf} n\'e
    \texttt{owl:hasValue}; le categorie (generi, piattaforme, modalit\`a) sono modellate
    come classi e non come insiemi chiusi di individui;
  \item $\mathcal{Q}$ (cardinalit\`a qualificate): le due restrizioni di cardinalit\`a
    non specificano \texttt{owl:onClass}.
\end{itemize}
Gli assiomi \texttt{owl:hasKey} non modificano l'espressivit\`a: sono \emph{DL-safe},
cio\`e si applicano solo agli individui nominati (Sezione~\ref{sec:haskey}).

\subsection{Profili OWL 2: perch\'e serve il reasoning a due livelli}

OWL 2 definisce tre profili trattabili (EL, QL, RL). L'ontologia non appartiene a
nessuno di essi:
\begin{itemize}[nosep]
  \item \textbf{OWL 2 EL} esclude ruoli inversi, complemento e cardinalit\`a;
  \item \textbf{OWL 2 QL} esclude, tra l'altro, le property chain e la transitivit\`a;
  \item \textbf{OWL 2 RL} ammette $\exists p.C$ solo come \emph{sottoclasse} e non ammette
    datatype restriction n\'e complemento in posizione di sottoclasse.
\end{itemize}

Il caso di OWL 2 RL \`e quello rilevante per il progetto. Un'equivalenza
$X \equiv \texttt{VideoGame} \sqcap \exists p.C$ contiene due inclusioni:
\begin{itemize}[nosep]
  \item $\texttt{VideoGame} \sqcap \exists p.C \sqsubseteq X$ (condizione sufficiente):
    ammessa in RL, \`e quella che serve per \emph{classificare} (regole
    \texttt{cls-int1} e \texttt{cls-svf1});
  \item $X \sqsubseteq \exists p.C$ (condizione necessaria): non ammessa in RL, perch\'e
    richiederebbe di introdurre individui anonimi.
\end{itemize}
Un reasoner RL come \texttt{owlrl} applica quindi correttamente la direzione utile
alla classificazione e ignora l'altra. Restano invece fuori dal profilo, in entrambe
le direzioni, le datatype restriction (\texttt{HighlyRatedGame}, \texttt{RecentGame})
e il complemento usato come condizione sufficiente (\texttt{StandaloneGame}): per
questi costrutti il progetto usa il secondo livello di reasoning basato su SPARQL
CONSTRUCT.

Le cardinalit\`a massime ${\le}1$ in posizione di superclasse, al contrario, sono
ammesse in RL: la regola \texttt{cls-maxc2} identifica (\texttt{owl:sameAs}) i due valori
di un individuo che ne avesse pi\`u di uno. Non producono per\`o classificazioni,
perch\'e \texttt{GameWithSingleDeveloper} ed \texttt{ExclusiveRelease} hanno sole
condizioni necessarie.

\subsection{Vincoli globali di OWL 2 DL}

Per garantire la decidibilit\`a, OWL 2 DL impone vincoli globali sull'uso delle
propriet\`a. Una propriet\`a \`e \emph{non semplice} se \`e transitiva, se \`e
super-propriet\`a di una catena, oppure se \`e inversa o super-propriet\`a di una
propriet\`a non semplice. Nello schema le propriet\`a non semplici sono sette:
\texttt{sequelOf}, \texttt{hasSequel}, \texttt{partOfFranchise}, \texttt{subsidiaryOf},
\texttt{sharedFranchiseWith}, \texttt{sharesDeveloperWith}, \texttt{sharesPublisherWith}.

I vincoli rispettati dallo schema sono:
\begin{enumerate}[nosep]
  \item \textbf{Propriet\`a semplici}: nessuna propriet\`a non semplice compare in
    assiomi di asimmetria, irriflessivit\`a, funzionalit\`a (diretta o inversa),
    cardinalit\`a o \texttt{hasKey}. Per questo le tre transitive e \texttt{hasSequel}
    non sono dichiarate asimmetriche n\'e irriflessive (Sezione~\ref{sec:tbox_proprieta});
  \item \textbf{Regolarit\`a delle catene}: le tre catene hanno la forma
    $R_1 \circ R_2 \sqsubseteq S$ con $S$ distinta da $R_1$ e $R_2$ e senza dipendenze
    cicliche tra catene, quindi la gerarchia di ruoli \`e regolare;
  \item \textbf{Nessuna InverseFunctionalProperty su datatype property}: OWL 2 DL non
    prevede l'inversa di una data property; l'identit\`a \`e espressa con \texttt{hasKey}.
\end{enumerate}

La verifica \`e stata eseguita con uno script basato su \texttt{rdflib}, che calcola
l'insieme delle propriet\`a non semplici e ne controlla l'uso negli assiomi elencati:
il risultato \`e di 0 violazioni.

\subsection{Limiti Noti}

Per l'uso con un reasoner DL completo (HermiT, Pellet) restano due aspetti:
\begin{itemize}[nosep]
  \item \textbf{Datatype \texttt{xsd:date}}. \`E il range di \texttt{releaseDate},
    \texttt{reviewDate} ed \texttt{eventDate} e la base della restrizione di
    \texttt{RecentGame}, ma non fa parte della \emph{datatype map} di OWL 2, che per gli
    istanti temporali prevede solo \texttt{xsd:dateTime} e \texttt{xsd:dateTimeStamp}.
    HermiT lo rifiuta come datatype non supportato; la soluzione \`e migrare schema e
    dati a \texttt{xsd:dateTime};
  \item \textbf{Entit\`a esterne non dichiarate}. Le entit\`a usate negli allineamenti
    (\texttt{wd:*}, \texttt{dcterms:*}) non hanno una dichiarazione di tipo esplicita;
    strumenti come l'OWL API le dichiarano implicitamente in fase di caricamento.
\end{itemize}

% Reasoning e interrogazione
\section{Reasoning}
\label{sec:reasoning}

\subsection{Architettura a Due Livelli}

Il sistema di reasoning adotta un'architettura a due livelli che combina un
reasoner OWL 2 RL standard con un reasoner CONSTRUCT personalizzato per i
costrutti non coperti dal profilo RL.

\subsubsection{Livello 1: owlrl (OWL 2 RL)}

\texttt{owlrl} \citep{owlrl} implementa la chiusura deduttiva per il profilo
OWL 2 RL \citep{w3c-owl2-rl}, materializzando le triple inferite direttamente
nel grafo RDF. I costrutti supportati includono:
\begin{itemize}[nosep]
  \item \texttt{rdfs:subClassOf} --- propagazione del tipo verso l'alto;
  \item \texttt{owl:someValuesFrom} --- inferenza da restrizioni esistenziali
    (es. se G \texttt{wonAward} A, allora G \texttt{rdf:type AwardWinningGame});
  \item \texttt{owl:hasValue} --- inferenza da valore specifico;
  \item \texttt{owl:intersectionOf} --- scomposizione di intersezioni;
  \item \texttt{owl:inverseOf} --- materializzazione di propriet\`a inverse;
  \item \texttt{owl:SymmetricProperty} e \texttt{owl:TransitiveProperty};
  \item \texttt{owl:propertyChainAxiom} --- catene di propriet\`a.
\end{itemize}

Sull'insieme di istanze di test (3 giochi + 2 developer + piattaforme), \texttt{owlrl} produce
un incremento da $\sim$997 a $\sim$2680 triple.

\subsubsection{Livello 2: construct\_reasoner.py (SPARQL CONSTRUCT)}

21 regole CONSTRUCT colmano le lacune di OWL 2 RL, operando su costrutti che
richiedono ragionamento al di fuori del profilo RL. La Tabella~\ref{tab:construct_rules}
elenca le regole.

\begin{longtable}{lp{6cm}l}
\caption{Regole CONSTRUCT del reasoner personalizzato}\label{tab:construct_rules}\\
\toprule
\textbf{\#} & \textbf{Descrizione} & \textbf{Cosa inferisce} \\
\midrule
\endfirsthead
\multicolumn{3}{c}{\tablename\ \thetable{} -- continua} \\
\toprule
\textbf{\#} & \textbf{Descrizione} & \textbf{Cosa inferisce} \\
\midrule
\endhead
\bottomrule
\endlastfoot

R1 & HighlyRatedGame da metacriticScore $\geq$ 85     & \texttt{?g a vg:HighlyRatedGame} \\
R2 & RecentGame da releaseDate $\geq$ 2020-01-01      & \texttt{?g a vg:RecentGame} \\
R3 & StandaloneGame: VideoGame senza sequelOf          & \texttt{?g a vg:StandaloneGame} (negation-as-failure) \\
R4 & SequelGame: VideoGame con sequelOf                & \texttt{?g a vg:SequelGame} \\
R5a--f& Classificazione per piattaforma (PC, Console, Mobile, Nintendo, PS, Xbox) & \texttt{?g a vg:PCGame}, ecc. \\
R6a--c& Classificazione per modalit\`a (Multiplayer, SinglePlayer, Coop) & \texttt{?g a vg:MultiplayerGame}, ecc. \\
R7 & Inverse hasSequel: se A sequelOf B $\rightarrow$ B hasSequel A & \texttt{?b vg:hasSequel ?a} \\
R8 & Inverse awardWonBy & \texttt{?a vg:awardWonBy ?g} \\
R9a--d& Inverse: genreOf, platformFor, engineUsedIn, modeInGame & completamento bidirezionale \\
R10 & sharedGenreWith (pairwise) & \texttt{?g1 vg:sharedGenreWith ?g2} \\
R11 & sharedEngineWith (pairwise) & \texttt{?g1 vg:sharedEngineWith ?g2} \\
\bottomrule
\end{longtable}

Solo R1--R3 coprono costrutti realmente esterni a OWL 2 RL. Le regole R4--R9
(classificazione per sequel, piattaforma e modalit\`a, completamento delle inverse)
ripetono inferenze che il livello 1 gi\`a produce: sono mantenute perch\'e il livello 2
possa essere eseguito anche da solo, senza dipendere dalla chiusura \texttt{owlrl}.
R10 e R11 aggiungono relazioni di comodo tra giochi con lo stesso genere o motore.

\subsection{Perch\'e Alcuni Costrutti Non Sono in OWL 2 RL}

Il profilo OWL 2 RL \`e deliberatamente limitato per garantire la trattabilit\`a
computazionale (polinomiale); la Sezione~\ref{sec:profilo_dl} ne discute i limiti in
dettaglio. Per le classi dell'ontologia la situazione \`e la seguente: due famiglie di
costrutti esulano dal profilo, mentre la terza vi rientra ma non produce classificazioni.

\begin{enumerate}[nosep]
  \item \textbf{Datatype restriction} (\texttt{HighlyRatedGame}, \texttt{RecentGame}).
    Le restrizioni su tipi di dato con \texttt{xsd:minInclusive} richiedono il
    confronto di valori literal, operazione non prevista dalle regole RL standard.
    Le regole R1 e R2 risolvono con un semplice filtro SPARQL (\texttt{FILTER(?s >= 85)});
  \item \textbf{complementOf puro} (\texttt{StandaloneGame}). OWL 2 RL supporta
    solo forme limitate di negazione. La regola R3 implementa la
    \textit{negation-as-failure} con \texttt{FILTER NOT EXISTS}, assumendo la
    completezza dei dati (closed-world assumption sulla propriet\`a \texttt{sequelOf});
  \item \textbf{maxCardinality} (\texttt{GameWithSingleDeveloper}, \texttt{ExclusiveRelease}).
    Una restrizione ${\le}1$ in posizione di superclasse \`e ammessa in OWL 2 RL
    (regola \texttt{cls-maxc2}: se un'istanza ha due valori, questi vengono identificati
    con \texttt{owl:sameAs}). Poich\'e le due classi hanno sole condizioni necessarie
    (\texttt{rdfs:subClassOf}), per\`o, nessun gioco viene classificato in esse: fungono
    da vincoli di integrit\`a, e nessuna regola CONSTRUCT \`e necessaria.
\end{enumerate}

\subsection{Esempio di Inferenza a Due Livelli}

L'insieme di istanze di test (Appendice~\ref{sec:app_abox}) permette di seguire il
ragionamento su un caso concreto. Il gioco \texttt{GameA} (Elden Ring) \`e asserito
come \texttt{VideoGame} con data di rilascio 2022-02-25, punteggio Metacritic 96,
sviluppatore \texttt{Dev1}, due modalit\`a (individui di tipo \texttt{SinglePlayerMode}
e \texttt{MultiplayerMode}), due piattaforme (\texttt{PS5}, di tipo
\texttt{PlayStationPlatform}, e \texttt{PC\_Windows}, di tipo \texttt{WindowsPlatform})
e un premio. Anche \texttt{GameC} \`e sviluppato da \texttt{Dev1}.
La Tabella~\ref{tab:inferenza} riporta le conclusioni, il motivo e il livello che le produce.

\begin{table}[h]
\centering
\small
\caption{Inferenze su \texttt{GameA}}
\label{tab:inferenza}
\begin{tabularx}{\textwidth}{l X c}
\toprule
\textbf{Conclusione} & \textbf{Motivo} & \textbf{Livello} \\
\midrule
\texttt{PlayStationGame}, \texttt{ConsoleGame} &
  \texttt{availableOn PS5}; \texttt{PS5} \`e \texttt{PlayStationPlatform} $\sqsubseteq$
  \texttt{ConsolePlatform}; definizioni con $\exists$\texttt{availableOn} & 1 \\
\texttt{PCGame} & \texttt{WindowsPlatform} $\sqsubseteq$ \texttt{PCPlatform} & 1 \\
\texttt{SinglePlayerGame}, \texttt{MultiplayerGame} & $\exists$\texttt{hasGameMode} verso le due modalit\`a & 1 \\
\texttt{AwardWinningGame} & $\exists$\texttt{wonAward.Award}; il premio \`e \texttt{IndustryAward} $\sqsubseteq$ \texttt{Award} & 1 \\
\texttt{HighlyRatedGame} & $96 \ge 85$ (datatype restriction) & 2 \\
\texttt{RecentGame} & $2022\text{-}02\text{-}25 \ge 2020\text{-}01\text{-}01$ (datatype restriction) & 2 \\
\texttt{StandaloneGame} & nessun \texttt{sequelOf} noto (negation-as-failure) & 2 \\
\midrule
\texttt{Dev1 developerOf GameA} & inversa di \texttt{developedBy} & 1 \\
\texttt{GameA sharesDeveloperWith GameC} & catena \texttt{developedBy} $\circ$ \texttt{developerOf} e simmetria & 1 \\
\bottomrule
\end{tabularx}
\end{table}

L'esempio mostra la divisione dei compiti: il livello 1 (\texttt{owlrl}) copre
gerarchie, restrizioni esistenziali, inverse e catene; il livello 2 (CONSTRUCT)
interviene solo dove servono confronti tra literal o negazione.

\subsection{Bug Critico: owl:sameAs Falsi e Soluzione}

\subsubsection{Il Problema}

Durante lo sviluppo \`e emerso un bug critico: \texttt{owlrl} inferiva
\texttt{owl:sameAs} falsi tra giochi distinti a causa degli assiomi di
identit\`a (\texttt{FunctionalProperty}, \texttt{InverseFunctionalProperty},
\texttt{hasKey}) dichiarati su \textbf{DatatypeProperty}.

\textbf{Meccanismo del bug}: nel profilo OWL 2 RL le uguaglianze tra individui
sono prodotte da tre regole:
\begin{itemize}[nosep]
  \item \texttt{prp-fp} (FunctionalProperty) --- se lo \emph{stesso soggetto} ha due
    valori per la propriet\`a, i due \emph{valori} vengono identificati;
  \item \texttt{prp-ifp} (InverseFunctionalProperty) --- se \emph{due soggetti}
    condividono un valore, i due \emph{soggetti} vengono identificati;
  \item \texttt{prp-key} (\texttt{hasKey}) --- due individui della classe con gli
    stessi valori di chiave vengono identificati.
\end{itemize}
\texttt{owlrl} tratta i literal come nodi del grafo alla stregua di individui.
Applicate a datatype property, \texttt{prp-ifp} e \texttt{prp-key} identificano
giochi distinti che condividono il valore di un identificativo, mentre \texttt{prp-fp}
produce uguaglianze tra literal diversi dello stesso gioco (es.\ due date di rilascio
regionali). Ogni \texttt{owl:sameAs} spurio propaga poi tipi e propriet\`a tra gli
individui unificati (regole \texttt{eq-rep-*}), falsando la classificazione.

\textbf{Esempio}: nel test delle istanze (Appendice~\ref{sec:app_abox}) due
\texttt{owl:sameAs} spuri collegavano \texttt{GameA} agli altri due giochi,
propagando loro l'asserzione \texttt{wonAward}: prima del fix tutti e tre i giochi
risultavano \texttt{AwardWinningGame}; dopo il fix solo \texttt{GameA}, l'unico con
\texttt{wonAward} asserito, viene classificato come tale.

\subsubsection{La Soluzione}

La soluzione adottata (implementata in \texttt{enrich\_owl.py}) si articola in
quattro passi:

\begingroup\sloppy
\begin{enumerate}[nosep]
  \item \textbf{Nessuna InverseFunctionalProperty su datatype property}. Oltre a
    innescare \texttt{prp-ifp}, questo costrutto non \`e ammesso in OWL 2 DL;
    l'identit\`a dei giochi \`e espressa esclusivamente tramite
    \texttt{owl:hasKey} (su \texttt{hasWikidataId} e \texttt{hasSteamAppId});
  \item \textbf{Sezione ``DL-only axioms'' separata}. I 16 assiomi rimanenti
    (13 \texttt{FunctionalProperty} su datatype property e 3 \texttt{owl:hasKey})
    sono dichiarati in una sezione dedicata in coda a \texttt{videogames.owl},
    con un commento esplicativo;
  \item \textbf{Strip automatico prima del reasoning}. La funzione
    \texttt{strip\_dl\_only\_axioms()} rimuove questi assiomi dalla copia di
    lavoro in memoria prima del reasoning, prevenendo le inferenze errate;
  \item \textbf{Restore opzionale}. Con il flag \texttt{-{}-keep-dl-axioms}, la
    funzione \texttt{restore\_dl\_only\_axioms()} ripristina gli assiomi nel file
    serializzato; per default l'output resta privo di assiomi DL-only e quindi
    sicuro per \texttt{owlrl}.
\end{enumerate}
\endgroup

Lo schema \texttt{videogames.owl} conserva sempre la sezione DL-only, preservandone
il valore dichiarativo per l'uso con reasoner DL completi (HermiT, Pellet).

\subsubsection{Impatto}

Dopo il fix: 0 \texttt{owl:sameAs} falsi su tutto l'insieme di istanze di test (2715 triple finali,
nessun gioco unificato erroneamente). Le 12 famiglie di regole Python dirette in
\texttt{enrich\_owl.py} (\texttt{\_materialise}) classificano correttamente tutti
i giochi senza produrre falsi positivi.

\subsection{Materialiser Mirato in enrich\_owl.py}

Oltre al reasoning owlrl + CONSTRUCT, \texttt{enrich\_owl.py} implementa un
materialiser Python diretto con 12 famiglie di regole che operano senza passare
per SPARQL, risultando pi\`u veloce di \texttt{owlrl} full su insiemi di istanze di grandi
dimensioni:

\begin{enumerate}[nosep]
  \item AwardWinningGame $\leftarrow$ game \texttt{wonAward} \_;
  \item FranchiseGame $\leftarrow$ game \texttt{belongsTo} \_;
  \item sharedFranchiseWith $\leftarrow$ belongsTo $\circ$ includes
  \item sharesDeveloperWith $\leftarrow$ developedBy $\circ$ developerOf
  \item sharesPublisherWith $\leftarrow$ publishedBy $\circ$ publisherOf
  \item MultiplayerGame / SinglePlayerGame / CoopGame $\leftarrow$ \texttt{hasGameMode}
  \item PCGame / ConsoleGame / MobileGame $\leftarrow$ \texttt{availableOn}
  \item NintendoGame / PlayStationGame / XboxGame $\leftarrow$ \texttt{availableOn}
  \item HighlyRatedGame $\leftarrow$ \texttt{metacriticScore} $\geq$ 85
  \item RecentGame $\leftarrow$ \texttt{releaseDate} $\geq$ 2020-01-01
  \item SequelGame $\leftarrow$ game \texttt{sequelOf} \_
  \item StandaloneGame $\leftarrow$ VideoGame senza \texttt{sequelOf}
\end{enumerate}

Le regole pairwise (3--5) sono protette dal cap \texttt{-{}-max-group} (default 50)
per evitare esplosioni combinatorie su franchise con migliaia di giochi.

\section{Interrogare l'Ontologia con SPARQL}
\label{sec:interrogazione}
\sloppy

Il valore degli assiomi si vede al momento dell'interrogazione: dopo il reasoning,
le conoscenze implicite sono triple esplicite del grafo e le query SPARQL possono
usarle direttamente, senza ricostruire la logica dello schema. Gli esempi seguenti
usano il prefisso \texttt{vg:} e mostrano quattro tipi di conoscenza derivata.

\subsection{Classi Inferite al Posto dei Filtri}

\begin{lstlisting}[style=sparql, caption={Giochi multiplayer per PlayStation con punteggio alto}, label={lst:q_classi}]
PREFIX vg: <http://www.videogame-ontology.org/ontology#>
SELECT ?name ?score WHERE {
  ?g a vg:HighlyRatedGame , vg:MultiplayerGame , vg:PlayStationGame ;
     vg:gameName ?name ;
     vg:metacriticScore ?score .
}
ORDER BY DESC(?score) LIMIT 10
\end{lstlisting}

Senza reasoning la stessa domanda richiederebbe un filtro sul punteggio
($\ge 85$), un join con le modalit\`a di gioco e un join con le piattaforme per
verificarne il tipo (ad esempio \texttt{PS5} \texttt{a} \texttt{PlayStationPlatform}). Con le classi definite (Sezione~\ref{sec:classi_definite})
la query si riduce a tre asserzioni di tipo.

\subsection{Relazioni Derivate da Property Chain}

\begin{lstlisting}[style=sparql, caption={Giochi dello stesso sviluppatore di Elden Ring}, label={lst:q_chain}]
PREFIX vg: <http://www.videogame-ontology.org/ontology#>
SELECT DISTINCT ?other WHERE {
  ?er vg:gameName "Elden Ring" ;
      vg:sharesDeveloperWith ?g .
  ?g  vg:gameName ?other .
  FILTER(?g != ?er)
}
\end{lstlisting}

\texttt{sharesDeveloperWith} non \`e mai asserita nei dati: deriva dalla catena
\texttt{developedBy} $\circ$ \texttt{developerOf} (Sezione~\ref{sec:chains}).
Il \texttt{FILTER} esclude la coppia riflessiva che la catena produce per ogni gioco.

\subsection{Transitivit\`a: Materializzazione o Property Path}

\begin{lstlisting}[style=sparql, caption={Tutti i predecessori di un gioco}, label={lst:q_trans}]
PREFIX vg: <http://www.videogame-ontology.org/ontology#>
SELECT ?pred WHERE {
  ?g vg:gameName "Dark Souls III" ;
     vg:sequelOf ?pred .      # con reasoning (chiusura materializzata)
  # senza reasoning:  ?g vg:sequelOf+ ?pred .
}
\end{lstlisting}

Poich\'e \texttt{sequelOf} \`e transitiva, dopo il reasoning il grafo contiene gi\`a
tutte le coppie (gioco, predecessore) a qualsiasi distanza. Lo stesso risultato si
pu\`o ottenere senza reasoning con il \emph{property path} SPARQL
\texttt{vg:sequelOf+}, che calcola la chiusura al momento della query. La differenza
\`e concettuale: nel primo caso la transitivit\`a \`e una propriet\`a del dominio
dichiarata nell'ontologia e vale per ogni applicazione che la usa; nel secondo \`e
una scelta di chi scrive la query.

\subsection{Navigazione Tramite Inverse}

\begin{lstlisting}[style=sparql, caption={Studi con pi\`u giochi premiati}, label={lst:q_inv}]
PREFIX vg: <http://www.videogame-ontology.org/ontology#>
SELECT ?studio (COUNT(DISTINCT ?g) AS ?n) WHERE {
  ?dev vg:developerOf ?g ;
       vg:developerName ?studio .
  ?g a vg:AwardWinningGame .
}
GROUP BY ?studio ORDER BY DESC(?n) LIMIT 10
\end{lstlisting}

I dati asseriscono solo \texttt{developedBy} (dal gioco allo studio): la query parte
dallo studio grazie all'inversa \texttt{developerOf} e usa la classe definita
\texttt{AwardWinningGame} invece di un join esplicito con i premi.

\subsection{Open World e Negazione}

Una domanda come \textless{}giochi senza sequel\textgreater{} mette in luce la differenza
tra l'ontologia e le query. In OWL vale la \textit{Open World Assumption}: l'assenza di
una tripla \texttt{sequelOf} non implica che il gioco non abbia predecessori, quindi un
reasoner DL non classificherebbe nessun gioco come \texttt{StandaloneGame} senza
ulteriori assiomi di chiusura. Il secondo livello di reasoning adotta invece
esplicitamente una \textit{Closed World Assumption} locale (\texttt{FILTER NOT EXISTS},
Sezione~\ref{sec:reasoning}): \texttt{StandaloneGame} va quindi letta come
\textless{}gioco per cui non \`e noto alcun predecessore\textgreater{}.

\fussy

% Contesto: popolamento e applicazione
\section{Popolamento dell'Ontologia}
\label{sec:popolamento}

\subsection{Pipeline \texttt{populate\_wikidata.py}}

La popolazione delle istanze avviene tramite lo script \texttt{populate\_wikidata.py},
che interroga l'endpoint SPARQL pubblico di Wikidata. La pipeline \`e strutturata
in quattro fasi principali.

\subsubsection{Query SPARQL verso Wikidata}

Vengono eseguite 12 tipologie di query SPARQL, ciascuna mirata a estrarre una
specifica categoria di dati:

\begin{enumerate}[nosep]
  \item \textbf{Core} --- nomi, date di rilascio, sviluppatore e publisher (P31, P577, P178, P123);
  \item \textbf{Generi} --- mapping genere via P136;
  \item \textbf{Piattaforme} --- piattaforme via P400;
  \item \textbf{Personaggi} --- personaggi presenti nel gioco via P674;
  \item \textbf{Franchise} --- serie di appartenenza via P179;
  \item \textbf{Game Mode} --- modalit\`a di gioco via P404;
  \item \textbf{Metacritic} --- punteggio via P444 con qualificatore P447 (Metacritic);
  \item \textbf{Engine} --- motore grafico via P408;
  \item \textbf{Paese} --- paese d'origine via P495;
  \item \textbf{Sito web} --- sito ufficiale via P856;
  \item \textbf{Premi} --- premi ricevuti via P166;
  \item \textbf{Descrizioni} --- descrizioni testuali via \texttt{schema:description}.
\end{enumerate}

Per evitare i timeout dell'endpoint pubblico Wikidata (limite $\sim$60 secondi),
le query sono partizionate per \textbf{anno $\times$ mese}: 17 anni (2010--2026)
$\times$ 12 mesi = \textbf{204 iterazioni} per ciascuna tipologia di query. Tra ogni richiesta viene inserito un
\texttt{polite delay} di 1 secondo. Il tempo totale di esecuzione \`e di circa
\textbf{5 ore}.

\subsubsection{Deduplicazione}

Wikidata pu\`o restituire la stessa entit\`a con URI diversi ma stesso nome
normalizzato (es. entit\`a omonime restituite da query diverse). La pipeline
applica una normalizzazione del nome (trim degli spazi e lowercase) e rimuove
\textbf{4.420 entit\`a duplicate}, preservando per ciascun nome l'URI con il
maggior numero di triple associate.

\subsubsection{OWL-RL Reasoning}

Dopo la popolazione iniziale ($\sim$1.019.589 triple grezze), viene applicata
la chiusura deduttiva OWL 2 RL tramite \texttt{owlrl}:
\begin{itemize}[nosep]
  \item Materializzazione propriet\`a inverse: per ogni \texttt{developedBy},
    viene aggiunto \texttt{developerOf};
  \item Catene subPropertyOf: \texttt{gameName} $\rightarrow$ \texttt{name};
  \item Inferenze da classi definite (\texttt{AwardWinningGame}, \texttt{FranchiseGame}).
\end{itemize}

Il grafo cresce da $\sim$1M a $\sim$1.4M triple.

\subsubsection{Pruning}

Lo step finale di pruning (\texttt{prune\_owl.py}) rimuove triple non necessarie
a runtime:
\begin{itemize}[nosep]
  \item \texttt{owl:sameAs} riflessivi (la chiusura OWL 2 RL genera un sameAs
    di ogni entit\`a con s\'e stessa);
  \item \texttt{gameDescription}: testi lunghi non usati nelle query;
  \item \texttt{officialWebsite}: URI non necessari per la ricerca semantica.
\end{itemize}

Complessivamente vengono rimosse $\sim$253k triple.
Risultato finale: $\sim$1.173.959 triple (dataset completo 2010--2026).

\subsection{Demo Subset 2020--2026}

Per rispettare i limiti di RAM del deploy su Render (512 MB), la demo live utilizza
un sottoinsieme dell'ontologia limitato ai giochi dal 2020 in poi (\texttt{prune\_years.py}):
\begin{itemize}[nosep]
  \item $\sim$745.400 triple;
  \item $\sim$68.700 giochi;
  \item RAM a runtime: $\sim$250 MB (pyoxigraph).
\end{itemize}

\subsection{Performance: rdflib vs pyoxigraph}

La Tabella~\ref{tab:perf} confronta le performance dei due store RDF utilizzati
nel progetto. Il passaggio da \texttt{rdflib} (Python puro) a \texttt{pyoxigraph}
(binding Rust) ha prodotto un miglioramento drastico.

\begin{table}[h]
\centering
\caption{Confronto performance rdflib vs pyoxigraph}
\label{tab:perf}
\begin{tabular}{lccc}
\toprule
\textbf{Metrica}              & \textbf{rdflib} & \textbf{pyoxigraph (completo)} & \textbf{pyoxigraph (demo)} \\
\midrule
RAM a runtime                 & $\sim$1.550 MB  & $\sim$389 MB                   & $\sim$250 MB               \\
Tempo di caricamento          & $\sim$41s       & $\sim$2.2s                     & $\sim$1.5s                 \\
Query SPARQL tipica           & 1--23s          & 0.08--0.37s                    & 0.05--0.25s                \\
Speedup query                 & ---             & 5--62$\times$                  & ---                        \\
\bottomrule
\end{tabular}
\end{table}

\texttt{pyoxigraph} \citep{pyoxigraph} garantisce un caricamento quasi istantaneo
($<$3 secondi contro 41 secondi) e query 5--62 volte pi\`u veloci, rendendo
l'esperienza utente fluida anche su hardware limitato.

\section{Architettura del Sistema}
\label{sec:architettura}

\subsection{Diagramma Architetturale}

L'architettura del sistema segue un classico pattern a tre livelli (three-tier):
frontend di presentazione, backend API REST e store ontologico. Il diagramma
seguente illustra il flusso dati:

\begin{center}
\begin{minipage}{0.9\textwidth}
\small
\begin{verbatim}
+----------------------------------------------------+
|  Frontend (React + Tailwind + react-force-graph)   |
|  [SearchBar]  [ResultList]   [Knowledge Graph 2D]  |
+--------------------------+-------------------------+
                           | HTTP REST
+--------------------------v-------------------------+
|               Backend (FastAPI)                    |
|  +---------------------------------------------+  |
|  |  SPARQL Agent (GPT-4.1-mini)                |  |
|  |  NL -> SPARQL -> Validate -> Execute        |  |
|  |  (retry automatico, max 3 tentativi)        |  |
|  +---------------------------------------------+  |
|  +------------------+  +-----------------------+  |
|  | OntologyService  |  |    GraphBuilder       |  |
|  |  (pyoxigraph)    |  | (nodi + archi grafo)  |  |
|  +------------------+  +-----------------------+  |
|  +------------------+  +-----------------------+  |
|  | Upstash Redis    |  | Wikipedia Image API   |  |
|  | (cache asincr.)  |  | (fallback cover art)  |  |
|  +------------------+  +-----------------------+  |
+--------------------------+-------------------------+
                           |
+--------------------------v-------------------------+
|  Ontologia (videogames_pruned_2020.owl)            |
|  Store: pyoxigraph (Rust)  ~745k triple            |
+----------------------------------------------------+
\end{verbatim}
\end{minipage}
\end{center}

\subsection{Stack Tecnologico}

La Tabella~\ref{tab:stack} riassume le tecnologie impiegate per ciascun livello.

\begin{table}[h]
\centering
\caption{Stack tecnologico del sistema}
\label{tab:stack}
\begin{tabular}{ll}
\toprule
\textbf{Layer}      & \textbf{Tecnologia}                         \\
\midrule
Frontend            & React 18, TypeScript, Vite, Tailwind CSS    \\
Grafo interattivo   & react-force-graph-2d                        \\
Backend             & FastAPI (Python 3.12+, async)               \\
Ontologia           & pyoxigraph (Rust), OWL 2, owlrl (rdflib)    \\
LLM / Agente        & OpenAI GPT-4.1-mini                         \\
Cache               & Upstash Redis (async) + cache in-memory      \\
Fonte dati          & Wikidata SPARQL                             \\
Deploy              & Render (backend) + Vercel (frontend)         \\
\bottomrule
\end{tabular}
\end{table}

\subsection{Flusso Dati}

Il flusso di una tipica interrogazione \`e il seguente:
\begin{enumerate}[nosep]
  \item L'utente digita una domanda in linguaggio naturale (NL) nella SearchBar del frontend;
  \item Il frontend invia la domanda al backend via HTTP POST su \texttt{/api/query};
  \item Lo \textbf{SPARQL Agent} (GPT-4.1-mini) genera una query SPARQL a partire dalla domanda NL, utilizzando il contesto dell'ontologia (classi, propriet\`a, esempi);
  \item Il \textbf{SPARQL Validator} controlla la correttezza sintattica della query generata; se invalida, l'agente ritenta (max 3 tentativi);
  \item La query viene eseguita su \textbf{pyoxigraph} contro il file OWL;
  \item Il \textbf{GraphBuilder} trasforma i risultati SPARQL in nodi e archi per il frontend (con metadati di tipo, label, colore);
  \item Il frontend riceve i risultati e li visualizza nella ResultList e nel Knowledge Graph 2D interattivo.
\end{enumerate}

\subsection{Cache a Due Livelli}

Per ottimizzare le performance, il sistema adotta una cache a due livelli:
\begin{itemize}[nosep]
  \item \textbf{Upstash Redis} (cloud, persistente, TTL 7 giorni): memorizza i risultati delle query SPARQL e le immagini (cover art) recuperate da Wikipedia. I risultati nulli non vengono cachati per consentire retry successivi;
  \item \textbf{Cache in-memory} (Python dict): mantiene in RAM label e tipi derivati dal grafo RDF per accessi rapidissimi durante la costruzione del grafo.
\end{itemize}

\subsection{Tabella dei Componenti}

La Tabella~\ref{tab:componenti} elenca i principali componenti software con le rispettive responsabilit\`a.

\begin{table}[h]
\centering
\caption{Componenti principali e responsabilit\`a}
\label{tab:componenti}
\begin{tabularx}{\textwidth}{lX}
\toprule
\textbf{Componente}          & \textbf{Responsabilit\`a}                                        \\
\midrule
\texttt{sparql\_agent.py}    & Traduzione NL $\rightarrow$ SPARQL via GPT-4.1-mini, retry       \\
\texttt{sparql\_validator.py}& Validazione sintattica SPARQL pre-esecuzione                     \\
\texttt{ontology\_service.py}& Caricamento OWL, esecuzione query via pyoxigraph                 \\
\texttt{graph\_builder.py}   & Costruzione nodi/archi da risultati RDF per il frontend          \\
\texttt{image\_cache.py}     & Recupero e cache immagini da Wikipedia API                       \\
\texttt{SearchBar.tsx}       & Input NL, ricerca con autocomplete                               \\
\texttt{ResultList.tsx}      & Visualizzazione risultati testuali con espansione dettaglio      \\
\texttt{KnowledgeGraph.tsx}  & Grafo 2D interattivo con force layout, zoom, filtri              \\
\texttt{NodeDetail.tsx}      & Pannello laterale con dettaglio nodo (propriet\`a, relazioni)    \\
\texttt{NodeContextMenu.tsx} & Menu contestuale (click destro): espandi, rimuovi nodo           \\
\texttt{SparqlViewer.tsx}    & Visualizzazione query SPARQL generata (espandibile)              \\
\bottomrule
\end{tabularx}
\end{table}

\subsection{API Endpoint}

La Tabella~\ref{tab:endpoint} riassume gli endpoint REST esposti dal backend.

\begin{table}[h]
\centering
\caption{Endpoint API}
\label{tab:endpoint}
\begin{tabular}{lll}
\toprule
\textbf{Metodo} & \textbf{Endpoint}             & \textbf{Descrizione}                     \\
\midrule
POST   & \texttt{/api/query}         & Ricerca NL $\rightarrow$ risultati + grafo \\
GET    & \texttt{/api/node/\{uri\}}  & Dettaglio nodo + sottografo               \\
GET    & \texttt{/api/image-search}  & Cerca cover art da Wikipedia              \\
GET    & \texttt{/api/stats}         & Statistiche ontologia                     \\
GET    & \texttt{/api/health}        & Health check                              \\
\bottomrule
\end{tabular}
\end{table}

\section{Frontend e Backend}
\label{sec:frontend_backend}

\subsection{Frontend: React + TypeScript + Vite}

L'interfaccia utente \`e un'applicazione single-page costruita con \textbf{React 18},
\textbf{TypeScript}, \textbf{Vite} e \textbf{Tailwind CSS} \citep{react}. Il
componente centrale \`e il \textbf{Knowledge Graph 2D}, realizzato con la libreria
\texttt{react-force-graph-2d}, che visualizza i risultati delle query come un
grafo interattivo physics-based.

\subsubsection{Componenti Principali}

I sei componenti React che compongono l'interfaccia sono:

\begin{itemize}[nosep]
  \item \textbf{SearchBar} --- input di ricerca con autocomplete. L'utente digita
    domande in linguaggio naturale e riceve suggerimenti;
  \item \textbf{ResultList} --- elenca i risultati testuali della query con
    metadati (nome, anno, punteggio). Ogni risultato \`e cliccabile per
    espandere il dettaglio;
  \item \textbf{KnowledgeGraph} --- il grafo interattivo 2D con force-directed
    layout. Supporta zoom, pan, hover (evidenzia le connessioni dirette),
    e un pulsante ``recentra'' per zoom-to-fit;
  \item \textbf{NodeDetail} --- pannello laterale che mostra tutte le propriet\`a
    e relazioni del nodo selezionato. Include link cliccabili per espandere
    altri nodi nel grafo;
  \item \textbf{NodeContextMenu} --- menu contestuale al click destro su un nodo:
    opzioni per espandere il nodo (aggiungere i vicini al grafo) o rimuoverlo;
  \item \textbf{SparqlViewer} --- pannello espandibile che mostra la query SPARQL
    generata dall'agente, utile per debugging e trasparenza.
\end{itemize}

\subsubsection{Caratteristiche del Grafo Interattivo}

Il Knowledge Graph offre un'esperienza di navigazione ricca:

\begin{itemize}[nosep]
  \item \textbf{Nodi colorati per tipo}: VideoGame (blu), Developer (verde),
    Publisher (arancione), Genre (viola), Platform (rosso), ecc.;
  \item \textbf{VideoGame come rettangoli portrait}: i nodi gioco sono visualizzati
    come riquadri verticali con immagine di copertina (da Wikipedia) o emoji
    placeholder, differenziandosi dagli altri nodi circolari;
  \item \textbf{Animazione physics-based}: forza di attrazione/repulsione tra i
    nodi con smorzamento progressivo;
  \item \textbf{Long-press su mobile}: supporto touch per dispositivi mobili;
  \item \textbf{Filtro per tipo}: legenda interattiva che permette di
    nascondere/mostrare intere categorie di nodi con un click;
  \item \textbf{Hover}: evidenzia le connessioni dirette del nodo, portandolo
    in primo piano (z-order).
\end{itemize}

\subsection{Backend: FastAPI Async}

Il backend \`e sviluppato con \textbf{FastAPI} \citep{fastapi} in modalit\`a
asincrona, esponendo 6 endpoint REST (cfr.~Tabella~\ref{tab:endpoint},
Sezione~\ref{sec:architettura}). I componenti chiave sono:

\subsubsection{SPARQL Agent}

Il cuore intelligente del sistema \`e lo \texttt{sparql\_agent.py}, che utilizza
\textbf{GPT-4.1-mini} di OpenAI \citep{openai} per tradurre domande in linguaggio
naturale in query SPARQL. Il prompt di sistema include:
\begin{itemize}[nosep]
  \item Il contesto completo dell'ontologia (classi, propriet\`a, esempi);
  \item I prefissi SPARQL (\texttt{vg:}, \texttt{rdf:}, \texttt{rdfs:}, \texttt{xsd:});
  \item Esempi di domande e query corrispondenti (\textit{few-shot prompting});
  \item Istruzioni per generare solo la query SPARQL, senza testo aggiuntivo.
\end{itemize}

Il flusso include un meccanismo di \textbf{retry automatico} (max 3 tentativi):
\begin{enumerate}[nosep]
  \item L'agente genera una query SPARQL;
  \item Il \texttt{sparql\_validator} ne verifica la correttezza sintattica;
  \item Se la query \`e sintatticamente invalida o produce zero risultati,
    l'agente riceve il messaggio di errore e ritenta con un prompt correttivo.
\end{enumerate}

\subsubsection{GraphBuilder}

Il \texttt{graph\_builder.py} trasforma i risultati SPARQL (binding di variabili)
in una struttura JSON adatta al frontend:
\begin{itemize}[nosep]
  \item \textbf{Nodi}: URI, label, tipo (classe OWL), colore, forma, immagine;
  \item \textbf{Archi}: soggetto, predicato, oggetto, label della relazione.
\end{itemize}

Il builder arricchisce i nodi con metadati derivati dal grafo RDF (tipo tramite
\texttt{rdf:type}, label tramite le propriet\`a \texttt{*Name}).

\subsubsection{Cache a Due Livelli}

\begin{itemize}[nosep]
  \item \textbf{Upstash Redis} \citep{upstash}: cache cloud persistente con TTL
    di 7 giorni per risultati query e URL immagini;
  \item \textbf{Cache in-memory}: dizionario Python per label e tipi, popolato
    al caricamento dell'ontologia e consultabile in O(1).
\end{itemize}

\subsection{Endpoint API}

Si rimanda alla Tabella~\ref{tab:endpoint} nella Sezione~\ref{sec:architettura}
per l'elenco completo degli endpoint REST.

\section{Conclusione}
\label{sec:conclusione}

\subsection{Riepilogo dei Risultati}

Il progetto \textbf{Videogame Semantic Search} ha prodotto un sistema completo
di ricerca semantica che integra:

\begin{itemize}[nosep]
  \item \textbf{Un'ontologia OWL 2} espressiva e ben strutturata, evoluta da 11 a
    78 classi (61 primitive + 17 definite), con 31 object property, 26 datatype
    property, 13 coppie inverse, 3 property chain, 2 propriet\`a transitive,
    12 blocchi di disjointness e allineamento esterno a Wikidata e Dublin Core;
  \item \textbf{Un sistema di reasoning a due livelli} (OWL 2 RL + 21 regole
    CONSTRUCT) funzionante end-to-end, con 0 falsi \texttt{owl:sameAs} dopo
    il fix del bug critico degli assiomi DL-only;
  \item \textbf{Un knowledge graph popolato} da Wikidata con oltre 1.1 milioni
    di triple (dataset completo 2010--2026) e $\sim$745k triple nel sottoinsieme
    demo (2020--2026, $\sim$68.700 giochi);
  \item \textbf{Un'interfaccia web interattiva} con ricerca in linguaggio naturale,
    agente SPARQL intelligente, grafo 2D navigabile e filtri per tipo;
  \item \textbf{Performance ottimizzate} grazie al passaggio da rdflib a pyoxigraph
    (5--62$\times$ speedup sulle query, caricamento in $<$3s, RAM $\sim$250 MB).
\end{itemize}

\subsection{Possibili Estensioni Future}

Il sistema offre numerose direzioni di sviluppo futuro:

\begin{enumerate}[nosep]
  \item \textbf{Reasoner DL completo} --- integrare HermiT o Pellet tramite
    \texttt{owlready2} per ragionare sui costrutti DL-only (datatype restriction,
    complementOf, maxCardinality) senza bisogno delle regole CONSTRUCT
    compensate. Questo abiliterebbe anche la rilevazione automatica di
    inconsistenze tra le istanze;
  \item \textbf{Validazione delle istanze con SHACL} --- definire SHACL shapes per
    validare la conformit\`a delle istanze (es. ogni VideoGame deve avere
    almeno un \texttt{gameName}, le date devono essere nel formato corretto);
  \item \textbf{SWRL rules reali} --- sostituire le 21 regole CONSTRUCT con
    autentiche SWRL rules, sfruttando un reasoner che le supporti nativamente;
  \item \textbf{Estensione della pipeline Wikidata} --- aggiungere query per
    recuperare \texttt{hasSteamAppId}, \texttt{budgetTier}, \texttt{playerCount}
    e la relazione \texttt{sequelOf}, attualmente non popolate dai dati Wikidata;
  \item \textbf{Internazionalizzazione del frontend} --- supporto multilingua
    per l'interfaccia e le query in linguaggio naturale (inglese, francese, \dots{});
  \item \textbf{Visualizzazione avanzata delle classi definite} --- nel grafo
    interattivo, evidenziare con badge o colori speciali i nodi che sono stati
    classificati automaticamente dal reasoner (es. bordino dorato per
    \texttt{HighlyRatedGame}, icona controller per \texttt{MultiplayerGame});
  \item \textbf{SPARQL endpoint pubblico} --- esporre un endpoint SPARQL
    standard (RESTful o tramite Triplestore) per consentire query federate
    da applicazioni esterne.
\end{enumerate}

Il progetto dimostra come le tecnologie del Semantic Web --- ontologie formali,
reasoning automatico, knowledge graph e query SPARQL --- possano essere integrate
in un'applicazione web moderna e performante, offrendo un'esperienza di ricerca
semantica che va ben oltre la semplice ricerca per parole chiave.


% ── Appendici ──
\appendix
\section{Esempio di Istanze di Test}
\label{sec:app_abox}

Per la verifica end-to-end del sistema di reasoning, \`e stato utilizzato un
insieme di istanze di test minimale (\texttt{test\_abox.ttl}) contenente 3 giochi, 2 developer
e alcune piattaforme. Il seguente estratto in Turtle illustra la struttura:

\begin{lstlisting}[style=turtle, caption={Estratto del test delle istanze (Turtle)}, label={lst:testabox}]
@prefix vg: <http://www.videogame-ontology.org/ontology#> .

vg:GameA  a vg:VideoGame ;
    vg:gameName "Elden Ring" ;
    vg:releaseDate "2022-02-25"^^xsd:date ;
    vg:metacriticScore 96 ;
    vg:developedBy vg:Dev1 ;
    vg:hasGenre vg:RPGGenreInstance ;
    vg:hasGameMode vg:SinglePlayer, vg:OnlineMultiplayer ;
    vg:availableOn vg:PS5, vg:PC_Windows ;
    vg:wonAward vg:TGAAward2022 .

vg:GameB  a vg:VideoGame ;
    vg:gameName "Hollow Knight" ;
    vg:releaseDate "2017-02-24"^^xsd:date ;
    vg:metacriticScore 87 ;
    vg:developedBy vg:Dev2 ;
    vg:availableOn vg:NintendoSwitch ;
    vg:hasGameMode vg:SinglePlayer .

vg:GameC  a vg:VideoGame ;
    vg:gameName "Dark Souls III" ;
    vg:sequelOf vg:GameX ;
    vg:developedBy vg:Dev1 .

# individui tipizzati (necessari per le restrizioni esistenziali)
vg:PS5               a vg:PlayStationPlatform .
vg:PC_Windows        a vg:WindowsPlatform .
vg:NintendoSwitch    a vg:NintendoPlatform .
vg:SinglePlayer      a vg:SinglePlayerMode .
vg:OnlineMultiplayer a vg:MultiplayerMode .
vg:TGAAward2022      a vg:IndustryAward .
\end{lstlisting}

\subsection{Risultati della Verifica}

Dopo il reasoning completo (owlrl + CONSTRUCT), il grafo di test produce:

\begin{itemize}[nosep]
  \item \textbf{2.715 triple finali} (dalle $\sim$997 iniziali);
  \item \textbf{0 \texttt{owl:sameAs} falsi} --- il fix degli assiomi DL-only
    previene l'unificazione errata dei giochi;
  \item \textbf{Classificazioni corrette}:
    \begin{itemize}[nosep]
      \item \texttt{GameA} $\rightarrow$ MultiplayerGame, SinglePlayerGame,
        ConsoleGame, PlayStationGame, PCGame, HighlyRatedGame, RecentGame,
        AwardWinningGame, StandaloneGame;
      \item \texttt{GameB} $\rightarrow$ SinglePlayerGame, ConsoleGame,
        NintendoGame, HighlyRatedGame, StandaloneGame (non RecentGame: 2017);
      \item \texttt{GameC} $\rightarrow$ SequelGame (nessuna classificazione per
        piattaforma o modalit\`a: non sono asseriti \texttt{availableOn} n\'e
        \texttt{hasGameMode}).
    \end{itemize}
  \item \textbf{Propriet\`a derivate}:
    \begin{itemize}[nosep]
      \item \texttt{Dev1} \texttt{developerOf} \texttt{GameA}, \texttt{GameC}
        (per l'inversa di \texttt{developedBy});
      \item \texttt{GameA} \texttt{sharesDeveloperWith} \texttt{GameC} e viceversa
        (catena \texttt{developedBy} $\circ$ \texttt{developerOf} e simmetria);
      \item \texttt{GameX} \texttt{hasSequel} \texttt{GameC} (inversa); \texttt{GameX}
        \`e inoltre classificato \texttt{VideoGame} per effetto del range di
        \texttt{sequelOf}.
    \end{itemize}
\end{itemize}

\section{Sintassi DL delle Classi Definite}
\label{sec:app_dl}

Di seguito si riporta la sintassi DL compatta di tutte le 17 classi definite
dell'ontologia. La notazione $\sqcap$ indica intersezione ($\mathtt{owl:intersectionOf}$),
$\exists$ indica restrizione esistenziale ($\mathtt{owl:someValuesFrom}$),
$\neg$ indica complemento ($\mathtt{owl:complementOf}$), $\leq$ indica
cardinalit\`a massima ($\mathtt{owl:maxCardinality}$).

\vspace{4pt}
\noindent
\begin{minipage}{\textwidth}
\small
\begin{align*}
\text{MultiplayerGame}      &\equiv \text{VideoGame} \sqcap \exists\text{hasGameMode.MultiplayerMode} \\
\text{SinglePlayerGame}     &\equiv \text{VideoGame} \sqcap \exists\text{hasGameMode.SinglePlayerMode} \\
\text{CoopGame}             &\equiv \text{VideoGame} \sqcap \exists\text{hasGameMode.CoopMode} \\
\text{PCGame}               &\equiv \text{VideoGame} \sqcap \exists\text{availableOn.PCPlatform} \\
\text{ConsoleGame}          &\equiv \text{VideoGame} \sqcap \exists\text{availableOn.ConsolePlatform} \\
\text{MobileGame}           &\equiv \text{VideoGame} \sqcap \exists\text{availableOn.MobilePlatform} \\
\text{NintendoGame}         &\equiv \text{VideoGame} \sqcap \exists\text{availableOn.NintendoPlatform} \\
\text{PlayStationGame}      &\equiv \text{VideoGame} \sqcap \exists\text{availableOn.PlayStationPlatform} \\
\text{XboxGame}             &\equiv \text{VideoGame} \sqcap \exists\text{availableOn.XboxPlatform} \\
\text{SequelGame}           &\equiv \text{VideoGame} \sqcap \exists\text{sequelOf.VideoGame} \\
\text{StandaloneGame}       &\equiv \text{VideoGame} \sqcap \neg\exists\text{sequelOf.VideoGame} \\
\text{HighlyRatedGame}      &\equiv \text{VideoGame} \sqcap \exists\text{metacriticScore.}\geq_{85} \\
\text{RecentGame}           &\equiv \text{VideoGame} \sqcap \exists\text{releaseDate.}\geq_{2020-01-01} \\
\text{AwardWinningGame}     &\equiv \text{VideoGame} \sqcap \exists\text{wonAward.Award} \\
\text{FranchiseGame}        &\equiv \text{VideoGame} \sqcap \exists\text{belongsTo.Franchise} \\[4pt]
\text{GameWithSingleDeveloper} &\sqsubseteq \text{VideoGame} \sqcap \mathopen{\leq}1\,\text{developedBy} \\
\text{ExclusiveRelease}     &\sqsubseteq \text{VideoGame} \sqcap \mathopen{\leq}1\,\text{availableOn}
\end{align*}
\end{minipage}

\vspace{8pt}
\noindent
Le prime 15 classi utilizzano $\equiv$ (condizioni necessarie e sufficienti,
\texttt{owl:equivalentClass}); le ultime due usano $\sqsubseteq$ (sole condizioni
necessarie, \texttt{rdfs:subClassOf}).

\section{Assiomi dello Schema in Turtle}
\label{sec:app_turtle}

Lo schema \`e memorizzato in RDF/XML (\texttt{videogames.owl}). Gli estratti seguenti
lo riportano in Turtle, pi\`u leggibile, per mostrare come la sintassi DL delle sezioni
precedenti diventa un insieme di triple RDF. Restrizioni, intersezioni e liste sono
rappresentate da nodi blank (\texttt{[ ... ]} e \texttt{( ... )}). Si omettono
\texttt{rdfs:label} e \texttt{rdfs:comment}.

\subsection*{Classe definita con restrizione esistenziale}

\begin{lstlisting}[style=turtle]
vg:PlayStationGame a owl:Class ;
    rdfs:subClassOf vg:VideoGame ;
    owl:equivalentClass [ a owl:Class ;
        owl:intersectionOf ( vg:VideoGame
            [ a owl:Restriction ;
              owl:onProperty vg:availableOn ;
              owl:someValuesFrom vg:PlayStationPlatform ] ) ] .
\end{lstlisting}

$\texttt{PlayStationGame} \equiv \texttt{VideoGame} \sqcap \exists\,\texttt{availableOn.PlayStationPlatform}$.
L'intersezione \`e una lista RDF di due elementi: la classe \texttt{VideoGame} e
una restrizione anonima.

\subsection*{Classe definita con datatype restriction}

\begin{lstlisting}[style=turtle]
vg:HighlyRatedGame a owl:Class ;
    rdfs:subClassOf vg:VideoGame ;
    owl:equivalentClass [ a owl:Class ;
        owl:intersectionOf ( vg:VideoGame
            [ a owl:Restriction ;
              owl:onProperty vg:metacriticScore ;
              owl:someValuesFrom [ a rdfs:Datatype ;
                  owl:onDatatype xsd:integer ;
                  owl:withRestrictions ( [ xsd:minInclusive 85 ] ) ] ] ) ] .
\end{lstlisting}

Il filler di \texttt{someValuesFrom} non \`e una classe ma un \emph{data range}:
il sottoinsieme di \texttt{xsd:integer} vincolato dal facet \texttt{xsd:minInclusive 85}.

\subsection*{Classe definita con complemento}

\begin{lstlisting}[style=turtle]
vg:StandaloneGame a owl:Class ;
    rdfs:subClassOf vg:VideoGame ;
    owl:equivalentClass [ a owl:Class ;
        owl:intersectionOf ( vg:VideoGame
            [ a owl:Class ;
              owl:complementOf [ a owl:Restriction ;
                  owl:onProperty vg:sequelOf ;
                  owl:someValuesFrom vg:VideoGame ] ] ) ] .
\end{lstlisting}

$\texttt{StandaloneGame} \equiv \texttt{VideoGame} \sqcap \lnot\exists\,\texttt{sequelOf.VideoGame}$.

\subsection*{Condizione solo necessaria con cardinalit\`a}

\begin{lstlisting}[style=turtle]
vg:GameWithSingleDeveloper a owl:Class ;
    rdfs:subClassOf vg:VideoGame ,
        [ a owl:Restriction ;
          owl:onProperty vg:developedBy ;
          owl:maxCardinality "1"^^xsd:nonNegativeInteger ] .
\end{lstlisting}

Due assiomi \texttt{rdfs:subClassOf} distinti, nessun \texttt{owl:equivalentClass}:
\[ \texttt{GameWithSingleDeveloper} \sqsubseteq \texttt{VideoGame} \sqcap {\le}1\,\texttt{developedBy} \]
L'assenza di \texttt{owl:onClass} rende la cardinalit\`a non qualificata.

\subsection*{Propriet\`a transitiva e inversa}

\begin{lstlisting}[style=turtle]
vg:sequelOf a owl:ObjectProperty , owl:TransitiveProperty ;
    rdfs:domain vg:VideoGame ;
    rdfs:range vg:VideoGame ;
    owl:inverseOf vg:hasSequel .

vg:hasSequel a owl:ObjectProperty ;
    rdfs:domain vg:VideoGame ;
    rdfs:range vg:VideoGame .
\end{lstlisting}

\subsection*{Property chain}

\begin{lstlisting}[style=turtle]
vg:sharesDeveloperWith a owl:ObjectProperty , owl:SymmetricProperty ;
    rdfs:domain vg:VideoGame ;
    rdfs:range vg:VideoGame ;
    owl:propertyChainAxiom ( vg:developedBy vg:developerOf ) .
\end{lstlisting}

La catena \`e una lista ordinata: l'ordine delle propriet\`a conta
($\texttt{developedBy} \circ \texttt{developerOf}$).

\subsection*{Disjointness e chiavi}

\begin{lstlisting}[style=turtle]
[] a owl:AllDisjointClasses ;
    owl:members ( vg:ConsolePlatform vg:PCPlatform vg:MobilePlatform ) .

vg:Review owl:hasKey ( vg:reviewedGame vg:reviewerName ) .
\end{lstlisting}

\texttt{owl:AllDisjointClasses} dichiara disgiunte a coppie tutte le classi della lista
con un solo assioma, invece di $n(n-1)/2$ assiomi \texttt{owl:disjointWith}.


% ── Bibliografia ──
\bibliographystyle{plainnat}
\bibliography{refs}

\end{document}
